% !TeX root = ../tesis.tex
%%%%%%%%%%%%%%%%%%%%%%%%%%%%%%%%%%%%%%%%%%%%%%%%%%%%%%%%%%%%%%%%%%%
% CAPÍTULO 5 — RESULTADOS DEL MODELO VFT
% Archivo índice: declara el capítulo e inputea las secciones.
% No agregar contenido directamente aquí.
%%%%%%%%%%%%%%%%%%%%%%%%%%%%%%%%%%%%%%%%%%%%%%%%%%%%%%%%%%%%%%%%%%%
\chapter{Resultados del Modelo VFT}
\label{ch:resultados-vft}

El presente capítulo expone los resultados obtenidos mediante la aplicación del
\textbf{Modelo VFT} (\textit{Vanishing Fig-Tree}) sobre la red de transporte
multimodal de la \zmvm. Los indicadores se presentan en tres fases progresivas
cuya estructura responde a sus dependencias computacionales: la Fase~1 puede
calcularse de forma autónoma a partir de la geometría de la red; la Fase~2
requiere los pesos dinámicos derivados de la Fase~1; y la Fase~3 exige que
las dos anteriores estén completas para ejecutar los algoritmos de caminos
mínimos y análisis de vulnerabilidad.

Las secciones marcadas con fondo \colorbox{azulRevision}{azul tenue}
corresponden a resultados que se encuentran sujetos a revisión.

% =============================================================================
%  SECCIONES DEL CAPÍTULO
% =============================================================================
% !TeX root = ../tesis.tex
%%%%%%%%%%%%%%%%%%%%%%%%%%%%%%%%%%%%%%%%%%%%%%%%%%%%%%%%%%%%%%%%%%%
% SECCIÓN 5.1 — LA RED METROPOLITANA COMO OBJETO DE ANÁLISIS
%%%%%%%%%%%%%%%%%%%%%%%%%%%%%%%%%%%%%%%%%%%%%%%%%%%%%%%%%%%%%%%%%%%
\section{La Red Metropolitana como Objeto de Análisis}
\label{sec:intro-resultados}

La construcción del grafo topológico de la \zmvm constituye el punto de partida
de todos los indicadores que se presentan en este capítulo. Para transformar la
geografía estática de las rutas de transporte en un modelo operativo y ruteable,
el \textbf{Modelo VFT} implementa una arquitectura de procesamiento en tres
capas progresivas, cuyo flujo se sintetiza en la Figura~\ref{fig:arq-capas}.

% --- Figura: arquitectura en 3 capas ---
\begin{figure}[H]
    \centering
    \begin{subfigure}[b]{0.32\textwidth}
        \centering
        \includegraphics[width=\textwidth]{Figures/Cap5/arq_capa1.png}
        \caption{Capa 1: Adquisición y procesamiento geométrico.}
        \label{fig:arq-capa1}
    \end{subfigure}
    \hfill
    \begin{subfigure}[b]{0.32\textwidth}
        \centering
        \includegraphics[width=\textwidth]{Figures/Cap5/arq_capa2.png}
        \caption{Capa 2: Modelado topológico y dinámico.}
        \label{fig:arq-capa2}
    \end{subfigure}
    \hfill
    \begin{subfigure}[b]{0.32\textwidth}
        \centering
        \includegraphics[width=\textwidth]{Figures/Cap5/arq_capa3.png}
        \caption{Capa 3: Evaluación y salida de indicadores.}
        \label{fig:arq-capa3}
    \end{subfigure}
    \caption[Arquitectura en tres capas del Modelo VFT]{Representación esquemática
    de las tres capas de procesamiento del Modelo VFT: adquisición geométrica,
    modelado topológico con impedancia temporal, y evaluación de indicadores.}
    \label{fig:arq-capas}
    \fuentefigura{Elaboración propia.}
\end{figure}

% --- Estadísticas del grafo ---
\subsection{Caracterización del Grafo Construido}
\label{subsec:estadisticas-grafo}

La ingesta de datos desde \texttt{Apimetro} proporcionó \textbf{22{,}766 entidades
espaciales} que, tras el proceso de \termino{Snapping Lógico Geográfico} mediante
un árbol \texttt{KDTree} con tolerancia de 50~m, se consolidaron en un grafo
dirigido con las siguientes dimensiones:

\begin{itemize}
    \item \textbf{11{,}115 nodos topológicos} (estaciones operativas, eliminando
    nodos fantasma y redundancias geométricas).
    \item \textbf{45{,}679 aristas}, de las cuales 25{,}197 corresponden a
    segmentos de servicio real y 20{,}482 a aristas de transbordo peatonal
    generadas algorítmicamente.
    \item \textbf{11 sistemas de transporte} representados: \textsc{Metro},
    Metrobús, Trolebús, Tren Ligero, Tren Suburbano, Cablebús, Mexicable,
    Mexibús, \textsc{rtp}, Corredores Concesionados y Pumabús.
    \item Distribución modal: \textbf{12\% infraestructura masiva estructurada}
    (Metro, MB, TL, CBB, SUB) y \textbf{88\% transporte de superficie}
    (RTP, CC, Trolebús, Mexibús, Mexicable).
\end{itemize}

La Figura~\ref{fig:paneles-grafo} ilustra las cuatro perspectivas de
visualización del grafo resultante: la infraestructura agregada, la conectividad
interestación, y la direccionalidad de ida y regreso.

% --- Figura: 4 paneles del grafo ---
\begin{figure}[H]
    \centering
    \begin{subfigure}[b]{0.48\textwidth}
        \centering
        \includegraphics[width=\textwidth]{Figures/Cap5/panel1_infraestructura.png}
        \caption{Panel 1: Infraestructura agregada de la red.}
        \label{fig:panel1}
    \end{subfigure}
    \hfill
    \begin{subfigure}[b]{0.48\textwidth}
        \centering
        \includegraphics[width=\textwidth]{Figures/Cap5/panel2_conectividad.png}
        \caption{Panel 2: Conectividad interestación (aristas).}
        \label{fig:panel2}
    \end{subfigure}

    \vspace{0.5em}

    \begin{subfigure}[b]{0.48\textwidth}
        \centering
        \includegraphics[width=\textwidth]{Figures/Cap5/panel3_ida.png}
        \caption{Panel 3: Direccionalidad de servicio (sentido ida).}
        \label{fig:panel3}
    \end{subfigure}
    \hfill
    \begin{subfigure}[b]{0.48\textwidth}
        \centering
        \includegraphics[width=\textwidth]{Figures/Cap5/panel4_regreso.png}
        \caption{Panel 4: Direccionalidad de servicio (sentido regreso).}
        \label{fig:panel4}
    \end{subfigure}
    \caption[Visualización del grafo topológico de la \zmvm]{Cuatro perspectivas
    de visualización del grafo dirigido construido sobre la red multimodal de la
    \zmvm. El proceso de Snapping Lógico redujo los nodos en aproximadamente
    un 90\% respecto a la representación geométrica cruda.}
    \label{fig:paneles-grafo}
    \fuentefigura{Elaboración propia con \texttt{VFTModel} (corrida 2026-04-18).}
\end{figure}

% --- Tabla resumen de indicadores ---
\subsection{Estructura de Indicadores y Dependencias por Fase}
\label{subsec:tabla-indicadores}

Los ocho indicadores que componen el análisis topológico se organizan en tres
fases cuya ejecución es secuencialmente dependiente, como se muestra en el
Cuadro~\ref{tab:resumen-indicadores-cap5}.

\begin{table}[H]
    \centering
    \caption[Resumen de indicadores del Modelo VFT por fase]{Resumen de los
    indicadores del Modelo VFT, su fase de cálculo, dependencias y estado
    de avance en la investigación.}
    \label{tab:resumen-indicadores-cap5}
    \begin{tabularx}{\textwidth}{|l|l|X|l|}
        \hline
        \textbf{Indicador} & \textbf{Fase} & \textbf{Dependencias} & \textbf{Estado} \\
        \hline
        Validación de Impedancia    & 1 & Grafo base + fallbacks            & Calculado \\
        Cobertura Espacial ($C$)    & 1 & Geometría de estaciones           & Calculado \\
        Fuerza Capilar ($FC$)       & 1 & Grafo base                        & Calculado \\
        Factor de Desviación ($DF$) & 1 & Grafo base + impedancia           & Calculado \\
        \hline
        Coef. de Fricción Vial ($CF$)  & 2 & Fase 1 completa                & Calculado \\
        Penalización por Transf. ($T_b$) & 2 & Fase 1 + $CF$              & Formulado \\
        Node Score ($VFT\_score$)   & 2 & Fase 1 + datos MV1/MV2            & Planificado \\
        \hline
        Tiempo de Viaje Prom. ($T$) & 3 & Fases 1 y 2 completas             & Pendiente \\
        Centralidad Interm. ($B$)   & 3 & Fases 1 y 2 + $T$                 & Pendiente \\
        Vulnerabilidad ($\Delta E$) & 3 & Fases 1, 2 y $B$                  & Pendiente \\
        \hline
    \end{tabularx}
    \fuentefigura{Elaboración propia.}
\end{table}

% !TeX root = ../tesis.tex
%%%%%%%%%%%%%%%%%%%%%%%%%%%%%%%%%%%%%%%%%%%%%%%%%%%%%%%%%%%%%%%%%%%
% SECCIÓN 5.2.1 — FASE 1: INDICADORES ESPACIALES Y DE COBERTURA
%
% Este archivo es el índice de la sección. El contenido está
% dividido en los archivos de la carpeta 5-2-1-Fase1/:
%
%   5-2-1-1-Impedancia.tex       ← Validación del Motor VFT
%   5-2-1-2-Cobertura.tex        ← Nivel de Cobertura Espacial (C)
%   5-2-1-3-FuerzaCapilar.tex    ← Fuerza Capilar (FC_H y FC)
%   5-2-1-4-FactorDesviacion.tex ← Factor de Desviación (DF)
%%%%%%%%%%%%%%%%%%%%%%%%%%%%%%%%%%%%%%%%%%%%%%%%%%%%%%%%%%%%%%%%%%%
\section{Fase 1: Indicadores Espaciales y de Cobertura}
\label{sec:fase1-indicadores}

La Fase~1 agrupa los indicadores que pueden calcularse a partir de la geometría
de la red sin requerir datos de demanda ni ponderaciones dinámicas. Su punto de
partida es la validación del motor de impedancia, que garantiza la coherencia
física de los pesos asignados antes de proceder al cálculo de cada indicador.

% !TeX root = ../../tesis.tex
% =============================================================================
%  5.2.1.1 VALIDACIÓN DEL MOTOR DE IMPEDANCIA
% =============================================================================
\subsection{Validación del Motor de Impedancia Temporal}
\label{subsec:validacion-impedancia}

Para verificar la consistencia del \textbf{Motor de Impedancia VFT}, se calculó
la velocidad efectiva implícita de cada sistema tras aplicar el Coeficiente de
Fricción Vial ($C_f$) y se contrastó con los rangos operativos nominales
documentados por los organismos operadores de la \zmvm. La velocidad efectiva
se obtiene a partir de la Ecuación~\eqref{eq:velocidad-efectiva}:

\begin{equation}
    V_{\text{efectiva}} = \frac{V_{\text{fallback}}}{C_f}
    \label{eq:velocidad-efectiva}
\end{equation}

donde $V_{\text{fallback}}$ representa la velocidad comercial de flujo libre del
sistema y $C_f$ es el coeficiente de fricción según el tipo de derecho de vía.

\colorbox{azulRevision}{El Cuadro~\ref{tab:validacion-impedancia} sintetiza los resultados} de la
corrida del modelo del 16 de abril de 2026, ejecutada sobre
\textbf{25{,}197 segmentos de servicio real}.

\begin{table}[H]
    \centering
    \caption[Validación de velocidades efectivas por sistema de transporte]{
    Contraste entre la velocidad implícita post-fricción calculada por el motor
    VFT y los rangos operativos nominales de cada sistema. Corrida: 2026-04-16.}
    \label{tab:validacion-impedancia}
    \begin{tabular}{|l|r|r|c|c|}
        \hline
        \textbf{Sistema} & \textbf{$V$ modelo (km/h)} & \textbf{$V$ implícita (km/h)} &
        \textbf{Rango esperado} & \textbf{Estado} \\
        \hline
        \colorbox{azulRevision}{Tren Suburbano}   & 65.0 & 65.0 & [55--75]  & \checkmark\ OK \\
        \colorbox{azulRevision}{Metro (STC)}      & 36.0 & 36.0 & [30--42]  & \checkmark\ OK \\
        \colorbox{azulRevision}{Tren Ligero}      & 22.0 & 22.0 & [18--26]  & \checkmark\ OK \\
        \colorbox{azulRevision}{Trolebús}         & 18.0 & 13.0 & [14--22]  & $\approx$ Límite \\
        \colorbox{azulRevision}{Cablebús}         & 20.0 & 17.4 & [16--24]  & \checkmark\ OK \\
        \colorbox{azulRevision}{Mexicable}        & 20.0 & 17.4 & [16--24]  & \checkmark\ OK \\
        \colorbox{azulRevision}{Metrobús}         & 16.3 & 14.1 & [13--20]  & \checkmark\ OK \\
        \colorbox{azulRevision}{Mexibús}          & 16.3 & 14.1 & [13--20]  & \checkmark\ OK \\
        \colorbox{azulRevision}{\textsc{rtp}}     & 14.0 &  8.0 & [10--18]  & $\approx$ Revisar \\
        \colorbox{azulRevision}{Corredores Conc.} & 11.0 &  8.0 & [8--14]   & \checkmark\ OK \\
        \colorbox{azulRevision}{Pumabús}          & 11.0 &  8.0 & [8--14]   & \checkmark\ OK \\
        \hline
    \end{tabular}
    \fuentefigura{Elaboración propia con \texttt{VFTModel} (corrida 2026-04-16).
    Fuentes de referencia: STC Metro \textsc{cdmx}, Metrobús \textsc{cdmx},
    \textsc{rtp} \textsc{cdmx}, \textsc{ste} Trolebús.}
\end{table}

La Figura~\ref{fig:impedancia-sistemas} presenta la distribución completa de
velocidades efectivas y tiempos por segmento para cada sistema.

\begin{figure}[H]
    \centering
    \includegraphics[width=\textwidth]{Figures/Cap5/fig01_impedancia.png}
    \caption[Caracterización de la impedancia temporal por sistema]{(a)~Velocidad
    efectiva post-fricción por sistema, coloreada según el tipo de derecho de vía.
    (b)~Distribución de tiempos de viaje por segmento interestación (segmentos
    menores a 5~min, sin valores atípicos). Corrida: 2026-04-16.}
    \label{fig:impedancia-sistemas}
    \fuentefigura{Elaboración propia con \texttt{VFTModel}.}
\end{figure}

Los resultados son coherentes con la metodología adoptada. El caso más
sensible es el del \textbf{Trolebús}, cuya velocidad efectiva (~13.0~km/h con
$C_f = 1.38$) se sitúa ligeramente por debajo del límite inferior del rango
operativo documentado (14--22~km/h). Esta diferencia de aproximadamente
1~km/h (7\%) es atribuible a que el tipo de derecho de vía \textit{compartido}
no captura la preferencia operativa parcial que el trolebús mantiene en ciertas
vialidades de la \textsc{cdmx}. Para fines de análisis estructural de red,
esta desviación se encuentra dentro de los márgenes de calibración aceptables.

% !TeX root = ../../tesis.tex
% =============================================================================
%  5.2.1.2 NIVEL DE COBERTURA ESPACIAL
% =============================================================================
\subsection{\bajoRevision{Nivel de Cobertura Espacial ($C$)}}
\label{subsec:cobertura-espacial}
\notaRevision{Los resultados de esta subsección están sujetos a revisión con el
tutor. Los valores presentados corresponden a la corrida del modelo del
18 de abril de 2026 con un isócrono de 800~m y pueden ajustarse tras
validación metodológica.}

El \textbf{Nivel de Cobertura Espacial} ($C$) evalúa la equidad territorial de
la red de transporte masivo al cuantificar qué proporción del área de la
demarcación se encuentra dentro de un radio de caminata aceptable respecto
a alguna estación. Se define mediante la siguiente expresión:

\begin{equation}
    C = \frac{A_{\text{cubierta}}}{A_{\text{total}}} \times 100
    \label{eq:cobertura-res}
\end{equation}

donde $A_{\text{cubierta}}$ es el área resultante de la unión de los
isócronos peatonales (\termino{buffers} de 800~m) de todas las estaciones del sistema,
y $A_{\text{total}}$ es la superficie geodésica total de la demarcación
político-administrativa analizada.

El Cuadro~\ref{tab:peores-cdmx} identifica las cinco alcaldías de la
\textsc{cdmx} con menor cobertura del sistema de transporte masivo,
evidenciando las zonas de mayor exclusión territorial.

\begin{table}[H]
    \centering
    \caption[Alcaldías con menor cobertura de transporte masivo (CDMX)]{
    Cinco alcaldías de la Ciudad de México con menor porcentaje de cobertura
    territorial del sistema de transporte masivo. Isócrono: 800~m.
    Corrida: 2026-04-18.}
    \label{tab:peores-cdmx}
    \begin{tabular}{|l|r|r|r|}
        \hline
        \textbf{Demarcación} & \textbf{Área total (km²)} &
        \textbf{Área cubierta (km²)} & \textbf{Cobertura (\%)} \\
        \hline
        Milpa Alta              & 298.01 &  35.0    &  11.75 \\
        La Magdalena Contreras  &  63.37 &  16.33   &  25.76 \\
        Tlalpan                 & 314.25 &  94.31   &  30.01 \\
        Tláhuac                 &  85.78 &  43.85   &  51.12 \\
        Cuajimalpa de Morelos   &  71.10 &  37.12   &  52.20 \\
        \hline
    \end{tabular}
    \fuentefigura{Elaboración propia con \texttt{VFTModel} y datos de
    \texttt{Apimetro}.}
\end{table}

\begin{table}[H]
    \centering
    \caption[Municipios del EDOMEX con menor cobertura de transporte masivo]{
    Cinco municipios del Estado de México con menor porcentaje de cobertura
    territorial del sistema de transporte masivo. Isócrono: 800~m.
    Corrida: 2026-04-18.}
    \label{tab:peores-edomex}
    \begin{tabular}{|l|r|r|r|}
        \hline
        \textbf{Demarcación} & \textbf{Área total (km²)} &
        \textbf{Área cubierta (km²)} & \textbf{Cobertura (\%)} \\
        \hline
        Tonanitla       &   8.52 & 0.00 & 0.00 \\
        Capulhuac       &  21.50 & 0.00 & 0.00 \\
        Calimaya        & 104.26 & 0.00 & 0.00 \\
        Donato Guerra   & 181.36 & 0.00 & 0.00 \\
        Ayapango        &  50.83 & 0.00 & 0.00 \\
        \hline
    \end{tabular}
    \fuentefigura{Elaboración propia con \texttt{VFTModel} y datos de
    \texttt{Apimetro}.}
\end{table}

% =============================================================================
%  Desglose de Cobertura por Sistema de Transporte
% =============================================================================

\subsubsection{Sistema de Transporte Colectivo Metro (METRO)}
\label{subsubsec:cobertura-metro}

El Sistema de Transporte Colectivo Metro constituye la columna vertebral de la
movilidad masiva en la Ciudad de México. Debido a su naturaleza de red ferroviaria
subterránea y elevada con derecho de vía exclusivo, sus estaciones actúan como
nodos estructurales de alta capacidad que articulan los flujos intermodales de la
metrópoli. El análisis de su cobertura espacial mediante isócronos peatonales de
800~m permite identificar las demarcaciones que se benefician directamente de este
sistema, así como aquellas periferias y zonas centrales que presentan una
desconexión crítica respecto a la red de transporte más masiva de la entidad.

El Cuadro~\ref{tab:cobertura-metro-cdmx} detalla las cinco alcaldías de la Ciudad de
México con los niveles más altos de cobertura territorial respecto a la red del
Metro, evidenciando la concentración de la infraestructura ferroviaria en el núcleo
urbano denso de la cuenca.

\begin{table}[H]
    \centering
    \caption[Alcaldías con mayor cobertura del STC Metro (CDMX)]{
    Cinco alcaldías de la Ciudad de México con mayor porcentaje de cobertura
    territorial para el Sistema de Transporte Colectivo Metro. Isócrono: 800~m.
    Corrida: 2026-04-18.}
    \label{tab:cobertura-metro-cdmx}
    \begin{tabular}{|l|r|r|r|}
        \hline
        \textbf{Demarcación} & \textbf{Área total (km²)} & \textbf{Área cubierta (km²)} & \textbf{Cobertura (\%)} \\
        \hline
        Cuauhtémoc          & 32.50 &  28.35 &  87.23 \\
        Benito Juárez       & 26.68 &  21.34 &  79.99 \\
        Venustiano Carranza & 33.84 &  24.39 &  72.77 \\
        Iztacalco           & 23.08 &  11.66 &  50.52 \\
        Azcapotzalco        & 33.50 &  13.80 &  41.18 \\
        \hline
    \end{tabular}
    \fuentefigura{Elaboración propia con \texttt{VFTModel} y datos de \texttt{Apimetro}.}
\end{table}

Por su parte, la distribución geográfica de esta cobertura se visualiza en la
Figura~\ref{fig:mapa-cobertura-metro}, donde se aprecia la concentración de la
infraestructura en la zona central de la cuenca y el drástico desvanecimiento del
servicio hacia las demarcaciones del sur y del poniente.

\begin{figure}[H]
    \centering
    \includegraphics[width=0.85\textwidth]{Figures/Cap5/Cobertura/Mapas/NB02_fig02_METRO_cobertura.png}
    \caption[Mapa de cobertura territorial del STC Metro en la CDMX]{Mapa del
    nivel de cobertura territorial del Sistema de Transporte Colectivo Metro
    por demarcación en la Ciudad de México, calculado a partir de \termino{buffers}
    de caminata peatonal de 800~m.}
    \label{fig:mapa-cobertura-metro}
    \fuentefigura{Elaboración propia con \texttt{VFTModel} utilizando visualización geoespacial.}
\end{figure}

\subsubsection{Metrobús (MB)}
\label{subsubsec:cobertura-mb}

El sistema Metrobús, que opera bajo el modelo de Autobuses de Tránsito Rápido
(\termino{Bus Rapid Transit}, BRT), se ha consolidado como la segunda red de
transporte masivo más importante de la Ciudad de México. Dada su mayor flexibilidad
constructiva frente a los sistemas ferroviarios pesados, su trazado ha logrado
expandirse por corredores viales estratégicos, complementando y alimentando a la
red del Metro. No obstante, al evaluar su cobertura mediante isócronos peatonales
de 800~m, persisten notables brechas territoriales, dejando a varias demarcaciones
periféricas con acceso limitado o nulo a esta red troncal.

El Cuadro~\ref{tab:cobertura-mb-cdmx} presenta las cinco alcaldías con mayor
proporción de su territorio cubierto por el Metrobús, destacando la concentración
del servicio en las zonas centro, norte y oriente de la ciudad, siguiendo los ejes
viales primarios.

\begin{table}[H]
    \centering
    \caption[Alcaldías con mayor cobertura de Metrobús (CDMX)]{
    Cinco alcaldías de la Ciudad de México con mayor porcentaje de cobertura
    territorial para el sistema Metrobús (MB). Isócrono: 800~m.
    Corrida: 2026-04-18.}
    \label{tab:cobertura-mb-cdmx}
    \begin{tabular}{|l|r|r|r|}
        \hline
        \textbf{Demarcación} & \textbf{Área total (km²)} & \textbf{Área cubierta (km²)} & \textbf{Cobertura (\%)} \\
        \hline
        Cuauhtémoc          & 32.50 &  26.74 &  82.27 \\
        Benito Juárez       & 26.68 &  16.21 &  60.74 \\
        Iztacalco           & 23.08 &  13.95 &  60.44 \\
        Venustiano Carranza & 33.84 &  16.10 &  47.58 \\
        Gustavo A. Madero   & 87.84 &  38.75 &  44.12 \\
        \hline
    \end{tabular}
    \fuentefigura{Elaboración propia con \texttt{VFTModel} y datos de \texttt{Apimetro}.}
\end{table}

Visualmente, el alcance espacial del sistema se ilustra en la
Figura~\ref{fig:mapa-cobertura-mb}. El mapa confirma cómo las líneas de BRT se
concentran fuertemente en las zonas centro, norte y oriente —siguiendo los ejes
viales primarios—, mientras que las barreras topográficas y de traza urbana
mantienen a las demarcaciones del surponiente con una cobertura residual o
inexistente.

\begin{figure}[H]
    \centering
    \includegraphics[width=0.85\textwidth]{Figures/Cap5/Cobertura/Mapas/NB02_fig02_MB_cobertura.png}
    \caption[Mapa de cobertura territorial del Metrobús en la CDMX]{Mapa del
    nivel de cobertura territorial del sistema Metrobús por demarcación en la
    Ciudad de México, calculado a partir de \termino{buffers} de caminata
    peatonal de 800~m.}
    \label{fig:mapa-cobertura-mb}
    \fuentefigura{Elaboración propia con \texttt{VFTModel} utilizando visualización geoespacial.}
\end{figure}

\subsubsection{Cablebús (CBB)}
\label{subsubsec:cobertura-cbb}

El Cablebús representa un paradigma distinto dentro de la red de transporte masivo.
Al ser un sistema de teleférico urbano, su diseño no busca abarcar grandes
extensiones del territorio mediante una red interconectada, sino resolver problemas
de accesibilidad muy puntuales en zonas de alta marginación y topografía compleja
(principalmente en la periferia). Por lo tanto, su nivel de cobertura espacial a
escala metropolitana es inherentemente focalizado, actuando como alimentador de
alta capacidad hacia las líneas del Metro.

El Cuadro~\ref{tab:cobertura-cbb-cdmx} presenta las cinco alcaldías con mayor
cobertura del Cablebús. Como es de esperarse por la naturaleza del sistema, la gran
mayoría de las demarcaciones de la Ciudad de México carecen por completo de esta
infraestructura, por lo que incluso los valores más altos son reducidos en términos
porcentuales absolutos.

\begin{table}[H]
    \centering
    \caption[Alcaldías con mayor cobertura de Cablebús (CDMX)]{
    Cinco alcaldías de la Ciudad de México con mayor porcentaje de cobertura
    territorial para el sistema Cablebús (CBB). Isócrono: 800~m.
    Corrida: 2026-04-18.}
    \label{tab:cobertura-cbb-cdmx}
    \begin{tabular}{|l|r|r|r|}
        \hline
        \textbf{Demarcación} & \textbf{Área total (km²)} & \textbf{Área cubierta (km²)} & \textbf{Cobertura (\%)} \\
        \hline
        Gustavo A. Madero   &  87.84 & 11.31 &  12.88 \\
        Miguel Hidalgo      &  46.39 &  5.50 &  11.85 \\
        Iztapalapa          & 113.07 & 12.74 &  11.27 \\
        Álvaro Obregón      &  95.82 &  4.38 &   4.57 \\
        Venustiano Carranza &  33.84 &  0.00 &   0.00 \\
        \hline
    \end{tabular}
    \fuentefigura{Elaboración propia con \texttt{VFTModel} y datos de \texttt{Apimetro}.}
\end{table}

La Figura~\ref{fig:mapa-cobertura-cbb} ilustra esta focalización espacial. El mapa
revela cómo las áreas de influencia del Cablebús se restringen a corredores muy
específicos (como en Gustavo A. Madero, Iztapalapa y Álvaro Obregón), contrastando
drásticamente con el vacío de infraestructura aérea en el resto de la cuenca.

\begin{figure}[H]
    \centering
    \includegraphics[width=0.85\textwidth]{Figures/Cap5/Cobertura/Mapas/NB02_fig02_CBB_cobertura.png}
    \caption[Mapa de cobertura territorial del Cablebús en la CDMX]{Mapa del
    nivel de cobertura territorial del sistema Cablebús por demarcación en la
    Ciudad de México, calculado a partir de \termino{buffers} de caminata
    peatonal de 800~m.}
    \label{fig:mapa-cobertura-cbb}
    \fuentefigura{Elaboración propia con \texttt{VFTModel} utilizando visualización geoespacial.}
\end{figure}

\subsubsection{Red de Transporte de Pasajeros (RTP)}
\label{subsubsec:cobertura-rtp}

La Red de Transporte de Pasajeros (\textsc{rtp}) desempeña un rol fundamental en
la equidad espacial del sistema de movilidad de la Ciudad de México. A diferencia
de los sistemas confinados, la \textsc{rtp} opera como un servicio de superficie
con vocación social y de alimentación, diseñado específicamente para penetrar en
zonas periféricas, de topografía compleja y de menor densidad, donde la
implementación de infraestructura pesada es inviable. Su extensión territorial es
notablemente mayor que la de cualquier otro sistema de la metrópoli.

El Cuadro~\ref{tab:cobertura-rtp-cdmx} presenta las cinco alcaldías con mayor
proporción de cobertura por parte de la \textsc{rtp}, confirmando la concentración
del servicio en el centro y norte de la ciudad, donde la densidad de rutas y paradas
alcanza sus niveles más altos.

\begin{table}[H]
    \centering
    \caption[Alcaldías con mayor cobertura de RTP (CDMX)]{
    Cinco alcaldías de la Ciudad de México con mayor porcentaje de cobertura
    territorial para la Red de Transporte de Pasajeros (RTP). Isócrono: 800~m.
    Corrida: 2026-04-18.}
    \label{tab:cobertura-rtp-cdmx}
    \begin{tabular}{|l|r|r|r|}
        \hline
        \textbf{Demarcación} & \textbf{Área total (km²)} & \textbf{Área cubierta (km²)} & \textbf{Cobertura (\%)} \\
        \hline
        Miguel Hidalgo      &  46.39 &  41.90 &  90.32 \\
        Cuauhtémoc          &  32.50 &  29.32 &  90.22 \\
        Azcapotzalco        &  33.50 &  30.08 &  89.81 \\
        Gustavo A. Madero   &  87.84 &  75.86 &  86.37 \\
        Iztapalapa          & 113.07 &  93.66 &  82.83 \\
        \hline
    \end{tabular}
    \fuentefigura{Elaboración propia con \texttt{VFTModel} y datos de \texttt{Apimetro}.}
\end{table}

La Figura~\ref{fig:mapa-cobertura-rtp} muestra la dispersión territorial de este
sistema. A diferencia del Metro o del Metrobús, el mapa de la \textsc{rtp} revela
una capilaridad mucho mayor hacia el sur y poniente de la ciudad, confirmando su
rol estratégico como conector de las zonas más alejadas de la cuenca con la red
troncal.

\begin{figure}[H]
    \centering
    \includegraphics[width=0.85\textwidth]{Figures/Cap5/Cobertura/Mapas/NB02_fig02_RTP_cobertura.png}
    \caption[Mapa de cobertura territorial de la RTP en la CDMX]{Mapa del
    nivel de cobertura territorial de la Red de Transporte de Pasajeros
    (\textsc{rtp}) por demarcación en la Ciudad de México, calculado a partir
    de \termino{buffers} de caminata peatonal de 800~m.}
    \label{fig:mapa-cobertura-rtp}
    \fuentefigura{Elaboración propia con \texttt{VFTModel} utilizando visualización geoespacial.}
\end{figure}

\subsubsection{Trolebús (TROLE)}
\label{subsubsec:cobertura-trole}

El Servicio de Transportes Eléctricos (\textsc{ste}), a través de su red de
Trolebús, constituye uno de los ejes históricos y fundamentales de la
electromovilidad de cero emisiones en la capital. Su red se distribuye
estratégicamente sobre los principales ejes viales, concentrando su mayor área
de influencia en las zonas centro, norte y oriente de la ciudad.

\textbf{Nota metodológica:} Es importante precisar que, al momento de la
ejecución de este análisis espacial (18 de abril de 2026), el Portal de Datos
Abiertos de la Ciudad de México aún no había liberado los trazados geoespaciales
ni el padrón de paradas oficiales correspondientes a la \textbf{Línea 13} del
Trolebús. Por consiguiente, dicha línea no se encuentra contabilizada en los
isócronos de cobertura presentados en esta sección.

El Cuadro~\ref{tab:cobertura-trole-cdmx} enumera las cinco alcaldías con mayor
cobertura territorial de este sistema, todas ellas ubicadas en la zona central y
oriente de la ciudad, donde se concentra la red eléctrica de superficie.

\begin{table}[H]
    \centering
    \caption[Alcaldías con mayor cobertura de Trolebús (CDMX)]{
    Cinco alcaldías de la Ciudad de México con mayor porcentaje de cobertura
    territorial para el sistema Trolebús (TROLE). Isócrono: 800~m.
    Corrida: 2026-04-18.}
    \label{tab:cobertura-trole-cdmx}
    \begin{tabular}{|l|r|r|r|}
        \hline
        \textbf{Demarcación} & \textbf{Área total (km²)} & \textbf{Área cubierta (km²)} & \textbf{Cobertura (\%)} \\
        \hline
        Cuauhtémoc    &  32.50 & 22.72 &  69.90 \\
        Benito Juárez &  26.68 & 16.50 &  61.85 \\
        Azcapotzalco  &  33.50 & 17.48 &  52.18 \\
        Iztacalco     &  23.08 & 11.93 &  51.69 \\
        Coyoacán      & 125.17 & 25.17 &  46.71 \\
        \hline
    \end{tabular}
    \fuentefigura{Elaboración propia con \texttt{VFTModel} y datos de \texttt{Apimetro}.}
\end{table}

Por su parte, la Figura~\ref{fig:mapa-cobertura-trole} ilustra la huella espacial
del Trolebús. El mapa visibiliza la fuerte concentración del servicio en alcaldías
como Benito Juárez, Cuauhtémoc, Iztacalco e Iztapalapa (impulsada por el Trolebús
Elevado), contrastando drásticamente con la carencia de infraestructura eléctrica
hacia las periferias montañosas.

\begin{figure}[H]
    \centering
    \includegraphics[width=0.85\textwidth]{Figures/Cap5/Cobertura/Mapas/NB02_fig02_TROLE_cobertura.png}
    \caption[Mapa de cobertura territorial del Trolebús en la CDMX]{Mapa del
    nivel de cobertura territorial del sistema Trolebús por demarcación en la
    Ciudad de México, calculado a partir de \termino{buffers} de caminata
    peatonal de 800~m.}
    \label{fig:mapa-cobertura-trole}
    \fuentefigura{Elaboración propia con \texttt{VFTModel} utilizando visualización geoespacial.}
\end{figure}

\subsubsection{Corredores Concesionados (CC)}
\label{subsubsec:cobertura-cc}

Los Corredores Concesionados constituyen la red de rutas de transporte de
superficie operadas por empresas privadas con concesión del Gobierno de la Ciudad
de México. Con base en su lógica de trazado sobre los ejes viales primarios y
secundarios de la ciudad consolidada, este sistema presenta uno de los niveles de
cobertura territorial más altos de la metrópoli, especialmente en las alcaldías de
mayor densidad urbana del centro y el oriente. Más que un sistema residual, los
Corredores Concesionados actúan como la red capilar que penetra en los intersticios
que los sistemas troncales no pueden servir directamente.

El Cuadro~\ref{tab:cobertura-cc-cdmx} presenta las cinco alcaldías con mayor
cobertura territorial de este sistema, donde se aprecia la saturación práctica del
servicio en el núcleo urbano denso de la capital.

\begin{table}[H]
    \centering
    \caption[Alcaldías con mayor cobertura de Corredores Concesionados (CDMX)]{
    Cinco alcaldías de la Ciudad de México con mayor porcentaje de cobertura
    territorial para el sistema de Corredores Concesionados (CC). Isócrono: 800~m.
    Corrida: 2026-04-18.}
    \label{tab:cobertura-cc-cdmx}
    \begin{tabular}{|l|r|r|r|}
        \hline
        \textbf{Demarcación} & \textbf{Área total (km²)} & \textbf{Área cubierta (km²)} & \textbf{Cobertura (\%)} \\
        \hline
        Cuauhtémoc          & 32.50 & 32.47 &  99.90 \\
        Benito Juárez       & 26.68 & 26.32 &  98.65 \\
        Iztacalco           & 23.08 & 22.25 &  96.42 \\
        Miguel Hidalgo      & 46.39 & 44.33 &  95.57 \\
        Venustiano Carranza & 33.84 & 27.77 &  82.06 \\
        \hline
    \end{tabular}
    \fuentefigura{Elaboración propia con \texttt{VFTModel} y datos de \texttt{Apimetro}.}
\end{table}

La Figura~\ref{fig:mapa-cobertura-cc} confirma la distribución omnipresente de
este sistema en la ciudad consolidada. El mapa ilustra cómo los Corredores
Concesionados recubren de manera prácticamente total las alcaldías centrales,
mientras que el servicio se va diluyendo hacia las periferias con suelo de
conservación o menor densidad.

\begin{figure}[H]
    \centering
    \includegraphics[width=0.85\textwidth]{Figures/Cap5/Cobertura/Mapas/NB02_fig02_CC_cobertura.png}
    \caption[Mapa de cobertura territorial de Corredores Concesionados en la CDMX]{
    Mapa del nivel de cobertura territorial del sistema de Corredores
    Concesionados por demarcación en la Ciudad de México, calculado a partir de
    \termino{buffers} de caminata peatonal de 800~m.}
    \label{fig:mapa-cobertura-cc}
    \fuentefigura{Elaboración propia con \texttt{VFTModel} utilizando visualización geoespacial.}
\end{figure}

\subsubsection{Tren Ligero (TL)}
\label{subsubsec:cobertura-tl}

El Tren Ligero es un sistema de transporte férreo de capacidad media operado por
el Servicio de Transportes Eléctricos (\textsc{ste}). Su trazado consta de una
única línea que conecta la estación terminal Tasqueña de la Línea 2 del Metro con
el centro de la alcaldía Xochimilco. Debido a esta configuración lineal y su
confinamiento exclusivo en el sur de la cuenca, su área de influencia geográfica
es altamente focalizada.

El Cuadro~\ref{tab:cobertura-tl-cdmx} presenta las alcaldías con mayor nivel de
cobertura por parte del Tren Ligero. Únicamente Coyoacán, Xochimilco y Tlalpan
registran acceso al sistema dentro del isócrono peatonal de 800~m, mientras que
el resto de las demarcaciones carece por completo de este servicio.

\begin{table}[H]
    \centering
    \caption[Alcaldías con mayor cobertura de Tren Ligero (CDMX)]{
    Cinco alcaldías de la Ciudad de México con mayor porcentaje de cobertura
    territorial para el sistema Tren Ligero (TL). Isócrono: 800~m.
    Corrida: 2026-04-18.}
    \label{tab:cobertura-tl-cdmx}
    \begin{tabular}{|l|r|r|r|}
        \hline
        \textbf{Demarcación} & \textbf{Área total (km²)} & \textbf{Área cubierta (km²)} & \textbf{Cobertura (\%)} \\
        \hline
        Coyoacán            &  53.88 & 9.31 &  17.28 \\
        Xochimilco          & 114.03 & 6.82 &   5.98 \\
        Tlalpan             & 314.25 & 4.07 &   1.30 \\
        Venustiano Carranza &  33.84 & 0.00 &   0.00 \\
        Tláhuac             &  85.78 & 0.00 &   0.00 \\
        \hline
    \end{tabular}
    \fuentefigura{Elaboración propia con \texttt{VFTModel} y datos de \texttt{Apimetro}.}
\end{table}

La Figura~\ref{fig:mapa-cobertura-tl} ilustra espacialmente esta restricción
territorial. El mapa confirma cómo el servicio es exclusivo del corredor sur,
actuando principalmente como alimentador de la red masiva central para los
habitantes de las alcaldías sureñas.

\begin{figure}[H]
    \centering
    \includegraphics[width=0.85\textwidth]{Figures/Cap5/Cobertura/Mapas/NB02_fig02_TL_cobertura.png}
    \caption[Mapa de cobertura territorial del Tren Ligero en la CDMX]{Mapa del
    nivel de cobertura territorial del sistema Tren Ligero por demarcación en
    la Ciudad de México, calculado a partir de \termino{buffers} de caminata
    peatonal de 800~m.}
    \label{fig:mapa-cobertura-tl}
    \fuentefigura{Elaboración propia con \texttt{VFTModel} utilizando visualización geoespacial.}
\end{figure}

\subsubsection{Mexibús (MXB)}
\label{subsubsec:cobertura-mxb}

El sistema Mexibús es el equivalente al Metrobús en el Estado de México: una red
de autobuses de tránsito rápido (\termino{Bus Rapid Transit}, BRT) que opera sobre
corredores viales del ámbito metropolitano conurbado. Su presencia dentro de los
límites político-administrativos de la Ciudad de México es marginal, circunscrita
únicamente a los puntos de cruce o de conexión interestatal, por lo que su
cobertura territorial en las alcaldías de la capital es forzosamente residual.

El Cuadro~\ref{tab:cobertura-mxb-cdmx} identifica las alcaldías de la Ciudad de
México con mayor cobertura territorial del sistema Mexibús. Los valores reducidos
—ninguno superior al 7\%— confirman que este sistema opera fundamentalmente fuera
de los límites de la entidad y que su incidencia en la cobertura intraurbana
capitalina es mínima.

\begin{table}[H]
    \centering
    \caption[Alcaldías con mayor cobertura de Mexibús (CDMX)]{
    Cinco alcaldías de la Ciudad de México con mayor porcentaje de cobertura
    territorial para el sistema Mexibús (MXB). Isócrono: 800~m.
    Corrida: 2026-04-18.}
    \label{tab:cobertura-mxb-cdmx}
    \begin{tabular}{|l|r|r|r|}
        \hline
        \textbf{Demarcación} & \textbf{Área total (km²)} & \textbf{Área cubierta (km²)} & \textbf{Cobertura (\%)} \\
        \hline
        Iztacalco           &  23.08 & 1.52 & 6.58 \\
        Venustiano Carranza &  33.84 & 1.90 & 5.62 \\
        Gustavo A. Madero   &  87.84 & 3.23 & 3.67 \\
        Xochimilco          & 114.03 & 0.00 & 0.00 \\
        Tláhuac             &  85.78 & 0.00 & 0.00 \\
        \hline
    \end{tabular}
    \fuentefigura{Elaboración propia con \texttt{VFTModel} y datos de \texttt{Apimetro}.}
\end{table}

La Figura~\ref{fig:mapa-cobertura-mxb} confirma esta presencia puntual del
sistema dentro de la capital. El mapa ilustra cómo la única área de influencia
del Mexibús en la Ciudad de México corresponde a los cruces fronterizos del
nororiente, sin penetrar en el tejido urbano interior.

\begin{figure}[H]
    \centering
    \includegraphics[width=0.85\textwidth]{Figures/Cap5/Cobertura/Mapas/NB02_fig02_MEXIBÚS_cobertura.png}
    \caption[Mapa de cobertura territorial de Mexibús en la CDMX]{Mapa del
    nivel de cobertura territorial del sistema Mexibús por demarcación en la
    Ciudad de México, calculado a partir de \termino{buffers} de caminata
    peatonal de 800~m.}
    \label{fig:mapa-cobertura-mxb}
    \fuentefigura{Elaboración propia con \texttt{VFTModel} utilizando visualización geoespacial.}
\end{figure}

\subsubsection{Mexicable (MEXICABLE)}
\label{subsubsec:cobertura-mexicable}

El Mexicable es un sistema de teleférico urbano que opera principalmente en el
Estado de México, diseñado para sortear las barreras topográficas de las zonas
altas y densamente pobladas de la periferia metropolitana. Su incursión en la
Ciudad de México es mínima y de carácter estrictamente interconectivo, fungiendo
como enlace alimentador hacia las redes troncales en la zona norte de la capital.

El Cuadro~\ref{tab:cobertura-mexicable-cdmx} presenta las alcaldías de la capital
con mayor nivel de cobertura para el Mexicable. Como es natural dada su jurisdicción
y trazado, únicamente la alcaldía Gustavo A. Madero registra un área de influencia
efectiva dentro del isócrono peatonal de 800~m, mientras que el resto de las
demarcaciones presentan cobertura nula.

\begin{table}[H]
    \centering
    \caption[Alcaldías con mayor cobertura de Mexicable (CDMX)]{
    Cinco alcaldías de la Ciudad de México con mayor porcentaje de cobertura
    territorial para el sistema Mexicable. Isócrono: 800~m.
    Corrida: 2026-04-18.}
    \label{tab:cobertura-mexicable-cdmx}
    \begin{tabular}{|l|r|r|r|}
        \hline
        \textbf{Demarcación} & \textbf{Área total (km²)} & \textbf{Área cubierta (km²)} & \textbf{Cobertura (\%)} \\
        \hline
        Gustavo A. Madero   &  87.84 & 2.80 & 3.19 \\
        Venustiano Carranza &  33.84 & 0.00 & 0.00 \\
        Xochimilco          & 114.03 & 0.00 & 0.00 \\
        Tláhuac             &  85.78 & 0.00 & 0.00 \\
        Tlalpan             & 314.25 & 0.00 & 0.00 \\
        \hline
    \end{tabular}
    \fuentefigura{Elaboración propia con \texttt{VFTModel} y datos de \texttt{Apimetro}.}
\end{table}

La Figura~\ref{fig:mapa-cobertura-mexicable} visualiza la focalización extrema de
este servicio. El mapa ilustra claramente cómo la infraestructura aérea del
Mexicable se restringe a un pequeño polígono en el extremo norte de la ciudad,
reiterando su naturaleza como transporte de penetración desde la zona conurbada.

\begin{figure}[H]
    \centering
    \includegraphics[width=0.85\textwidth]{Figures/Cap5/Cobertura/Mapas/NB02_fig02_MEXICABLE_cobertura.png}
    \caption[Mapa de cobertura territorial del Mexicable en la CDMX]{Mapa del
    nivel de cobertura territorial del sistema Mexicable por demarcación en la
    Ciudad de México, calculado a partir de \termino{buffers} de caminata
    peatonal de 800~m.}
    \label{fig:mapa-cobertura-mexicable}
    \fuentefigura{Elaboración propia con \texttt{VFTModel} utilizando visualización geoespacial.}
\end{figure}

\subsubsection{Tren Suburbano (SUB)}
\label{subsubsec:cobertura-sub}

El Tren Suburbano de la \zmvm{} es un sistema ferroviario regional de alta
capacidad que conecta la terminal Buenavista, en la alcaldía Cuauhtémoc, con el
municipio de Cuautitlán, en el Estado de México. Su trazado, mayoritariamente
confinado al corredor norte de la cuenca metropolitana, limita su área de
influencia directa dentro de la Ciudad de México a un corredor muy acotado. Para
las alcaldías del sur y oriente, este sistema resulta prácticamente inaccesible
sin una transferencia modal significativa.

El Cuadro~\ref{tab:cobertura-sub-cdmx} detalla las cinco alcaldías de la Ciudad de
México con mayor nivel de cobertura por parte del Tren Suburbano, evidenciando la
concentración casi exclusiva del servicio en las demarcaciones del norte y centro
de la capital.

\begin{table}[H]
    \centering
    \caption[Alcaldías con mayor cobertura de Tren Suburbano (CDMX)]{
    Cinco alcaldías de la Ciudad de México con mayor porcentaje de cobertura
    territorial para el sistema Tren Suburbano (SUB). Isócrono: 800~m.
    Corrida: 2026-04-18.}
    \label{tab:cobertura-sub-cdmx}
    \begin{tabular}{|l|r|r|r|}
        \hline
        \textbf{Demarcación} & \textbf{Área total (km²)} & \textbf{Área cubierta (km²)} & \textbf{Cobertura (\%)} \\
        \hline
        Cuauhtémoc          &  32.50 & 2.01 & 6.18 \\
        Azcapotzalco        &  33.50 & 2.01 & 5.99 \\
        Xochimilco          & 114.03 & 0.00 & 0.00 \\
        Venustiano Carranza &  33.84 & 0.00 & 0.00 \\
        Tláhuac             &  85.78 & 0.00 & 0.00 \\
        \hline
    \end{tabular}
    \fuentefigura{Elaboración propia con \texttt{VFTModel} y datos de \texttt{Apimetro}.}
\end{table}

La Figura~\ref{fig:mapa-cobertura-sub} confirma la huella territorial de este
sistema. El mapa evidencia cómo la cobertura del Tren Suburbano apenas alcanza la
alcaldía Cuauhtémoc y el extremo norte de Azcapotzalco, mientras que el resto del
territorio capitalino permanece fuera de su radio de influencia peatonal.

\begin{figure}[H]
    \centering
    \includegraphics[width=0.85\textwidth]{Figures/Cap5/Cobertura/Mapas/NB02_fig02_SUB_cobertura.png}
    \caption[Mapa de cobertura territorial del Tren Suburbano en la CDMX]{Mapa
    del nivel de cobertura territorial del sistema Tren Suburbano por demarcación
    en la Ciudad de México, calculado a partir de \termino{buffers} de caminata
    peatonal de 800~m.}
    \label{fig:mapa-cobertura-sub}
    \fuentefigura{Elaboración propia con \texttt{VFTModel} utilizando visualización geoespacial.}
\end{figure}

\subsubsection{Tren Interurbano (INTERURBANO)}
\label{subsubsec:cobertura-interurbano}

El Tren Interurbano ``El Insurgente'' es un sistema ferroviario de carácter
regional diseñado para conectar el Valle de Toluca con la Ciudad de México. Su
trazo dentro de la capital es incipiente y se concentra exclusivamente en el sector
poniente, sirviendo como puerta de entrada a la metrópoli. Debido a su vocación
interurbana, su impacto en la cobertura espacial a nivel de alcaldías es altamente
localizado.

El Cuadro~\ref{tab:cobertura-interurbano-cdmx} presenta las alcaldías con mayor
nivel de cobertura para este sistema. Como reflejo de su ruta de ingreso, únicamente
Cuajimalpa de Morelos y Álvaro Obregón presentan áreas de influencia dentro del
isócrono peatonal de 800~m, dejando al resto de la ciudad sin cobertura directa
por este medio.

\begin{table}[H]
    \centering
    \caption[Alcaldías con mayor cobertura de Tren Interurbano (CDMX)]{
    Cinco alcaldías de la Ciudad de México con mayor porcentaje de cobertura
    territorial para el sistema Tren Interurbano (INTERURBANO). Isócrono: 800~m.
    Corrida: 2026-04-18.}
    \label{tab:cobertura-interurbano-cdmx}
    \begin{tabular}{|l|r|r|r|}
        \hline
        \textbf{Demarcación} & \textbf{Área total (km²)} & \textbf{Área cubierta (km²)} & \textbf{Cobertura (\%)} \\
        \hline
        Cuajimalpa de Morelos &  71.10 & 1.59 & 2.23 \\
        Álvaro Obregón        &  95.82 & 0.42 & 0.44 \\
        Venustiano Carranza   &  33.84 & 0.00 & 0.00 \\
        Xochimilco            & 114.03 & 0.00 & 0.00 \\
        Tláhuac               &  85.78 & 0.00 & 0.00 \\
        \hline
    \end{tabular}
    \fuentefigura{Elaboración propia con \texttt{VFTModel} y datos de \texttt{Apimetro}.}
\end{table}

La Figura~\ref{fig:mapa-cobertura-interurbano} visualiza este alcance espacial
limitado. El mapa hace evidente la focalización de la infraestructura en el límite
occidental de la cuenca, evidenciando su papel como colector foráneo más que como
distribuidor urbano capilar.

\begin{figure}[H]
    \centering
    \includegraphics[width=0.85\textwidth]{Figures/Cap5/Cobertura/Mapas/NB02_fig02_INTERURBANO_cobertura.png}
    \caption[Mapa de cobertura territorial del Tren Interurbano en la CDMX]{Mapa
    del nivel de cobertura territorial del sistema Tren Interurbano por
    demarcación en la Ciudad de México, calculado a partir de \termino{buffers}
    de caminata peatonal de 800~m.}
    \label{fig:mapa-cobertura-interurbano}
    \fuentefigura{Elaboración propia con \texttt{VFTModel} utilizando visualización geoespacial.}
\end{figure}

\subsubsection{Pumabús (PB)}
\label{subsubsec:cobertura-pb}

El sistema Pumabús es un servicio de transporte interno operado por la Universidad
Nacional Autónoma de México (UNAM), de carácter público y gratuito, cuya operación
es exclusiva del campus de Ciudad Universitaria. No obstante, debido a su relevancia
modal y estructural en el sur de la ciudad, se encuentra registrado formalmente
dentro del Portal de Datos Abiertos de la Ciudad de México, lo que permite integrar
su infraestructura de paradas y circuitos en el análisis macro de la conectividad
metropolitana.

El Cuadro~\ref{tab:cobertura-pb-cdmx} identifica las cinco alcaldías de la Ciudad
de México con mayor cobertura territorial respecto a la red del Pumabús,
evidenciando su completa focalización en el núcleo universitario y su vecindad
inmediata.

\begin{table}[H]
    \centering
    \caption[Alcaldías con mayor cobertura de Pumabús (CDMX)]{
    Cinco alcaldías de la Ciudad de México con mayor porcentaje de cobertura
    territorial para el sistema Pumabús (PB). Isócrono: 800~m.
    Corrida: 2026-04-18.}
    \label{tab:cobertura-pb-cdmx}
    \begin{tabular}{|l|r|r|r|}
        \hline
        \textbf{Demarcación} & \textbf{Área total (km²)} & \textbf{Área cubierta (km²)} & \textbf{Cobertura (\%)} \\
        \hline
        Coyoacán            &  53.88 & 10.85 &  20.13 \\
        Álvaro Obregón      &  95.82 &  2.25 &   2.34 \\
        Tlalpan             & 314.25 &  0.25 &   0.08 \\
        Venustiano Carranza &  33.84 &  0.00 &   0.00 \\
        Xochimilco          & 114.03 &  0.00 &   0.00 \\
        \hline
    \end{tabular}
    \fuentefigura{Elaboración propia con \texttt{VFTModel} y datos de \texttt{Apimetro}.}
\end{table}

La configuración geoespacial de esta influencia se presenta en la
Figura~\ref{fig:mapa-cobertura-pb}. El mapa hace patente la delimitación del
servicio al polígono escolar de Ciudad Universitaria, extendiendo su área de
influencia peatonal de 800~m únicamente hacia las franjas fronterizas de las tres
alcaldías circundantes.

\begin{figure}[H]
    \centering
    \includegraphics[width=0.85\textwidth]{Figures/Cap5/Cobertura/Mapas/NB02_fig02_PUMABUS_cobertura.png}
    \caption[Mapa de cobertura territorial del Pumabús en la CDMX]{Mapa del
    nivel de cobertura territorial del sistema Pumabús por demarcación en la
    Ciudad de México, calculado a partir de \termino{buffers} de caminata
    peatonal de 800~m.}
    \label{fig:mapa-cobertura-pb}
    \fuentefigura{Elaboración propia con \texttt{VFTModel} utilizando visualización geoespacial.}
\end{figure}

% !TeX root = ../../tesis.tex
% =============================================================================
%  5.2.1.3 FUERZA CAPILAR Y NODOS DE TRANSFERENCIA
% =============================================================================
\subsection{\bajoRevision{Fuerza Capilar y Nodos de Transferencia ($FC$)}}
\label{subsec:fuerza-capilar}
\notaRevision{Los resultados de esta subsección están sujetos a revisión con el
comité tutor. Los valores presentados corresponden a la corrida del modelo del
18 de abril de 2026.}

La \textbf{Fuerza Capilar} ($FC$) cuantifica el nivel de convergencia de flujos
de transporte en cada nodo o agrupación de nodos del grafo. El indicador se
evalúa desde dos perspectivas complementarias:

\textbf{A. Fuerza Capilar Geográfica ($FC_H$) — Macro-Hubs:} agrupa
topológicamente estaciones físicamente colindantes (radio de 100~m) en un
subgrafo $H$ que consolida los flujos de sus estaciones constituyentes,
representando el fenómeno urbano del transbordo físico (CETRAMs):

\begin{equation}
    FC_{H} = \sum_{j \in H} \bigl(k_{\text{in}}(j) + k_{\text{out}}(j)\bigr)
    - \text{Aristas Internas}
    \label{eq:fc-geografica}
\end{equation}

El Cuadro~\ref{tab:fuerza-capilar-hubs} detalla el comportamiento numérico de los
cinco principales Macro-Hubs de transferencia identificados en la red con base en
esta formulación, mientras que la Figura~\ref{fig:barplot-macro-hubs} permite
contrastar las magnitudes relativas de la concentración de flujos en los diez
nodos más significativos de la metrópoli.

\begin{table}[H]
    \centering
    \caption[Top 5 Macro-Hubs con mayor Fuerza Capilar Geográfica ($FC_H$)]{
    Los cinco Macro-Hubs de transferencia (nodos agrupados topológicamente a
    menos de 100~m) con mayor Fuerza Capilar Geográfica en la red metropolitana.
    Corrida: 2026-04-18.}
    \label{tab:fuerza-capilar-hubs}
    \fontsize{8.5pt}{10.5pt}\selectfont
    \begin{tabularx}{\textwidth}{|c|X|c|c|c|c|}
        \hline
        \textbf{ID Hub} & \textbf{Principales Estaciones Agrupadas} & \textbf{Sistemas} & \textbf{Nodos} & \textbf{Conexiones} \textit{(In / Out)} & \textbf{$FC_H$} \\
        \hline
        2174 & Zona Hospitales Sur (San Fernando, Chimalcoyotl, Coapa, Hosp. Manuel Gea Gtz., INER, Urgencias Respiratorias) & CC, RTP & 21 & 66 / 64 & \textbf{130} \\
        \hline
        136 & Eje Huipulco -- Estadio Azteca (Acoxpa, Huipulco, Estadio Azteca, Renato Leduc, San Alejandro) & CC, RTP, TL & 61 & 61 / 50 & \textbf{111} \\
        \hline
        1597 & Corredor Tlalpan Sur (Calle La Rosa, Tezoquipa, Galeana, Jojutla) & CC, RTP & 10 & 49 / 51 & \textbf{100} \\
        \hline
        259 & Nodo Barranca del Muerto (Barranca del Muerto, Alpes, Periférico Sur, José María Velasco, Miguel Ocaranza) & CC, RTP & 21 & 50 / 42 & \textbf{92} \\
        \hline
        129 & CETRAM Chapultepec (Metro Chapultepec, Acapulco, Agustín Melgar, Sonora, Veracruz, Tampico) & CC, METRO, RTP, TROLE & 29 & 37 / 40 & \textbf{77} \\
        \hline
    \end{tabularx}
    \vspace{1ex}
    \fuentefigura{Elaboración propia con \texttt{VFTModel} y datos de \texttt{Apimetro}.}
\end{table}

\begin{figure}[H]
    \centering
    \includegraphics[width=0.85\textwidth]{Figures/Cap5/FuerzaCapilar/Tabla/NB03_fig01_top10_macrohubs_barplot.png}
    \caption[Distribución de los principales Macro-Hubs por Fuerza Capilar]{
    Gráfico de barras comparativo de los diez principales Macro-Hubs de la red
    metropolitana, ordenados con base en su nivel acumulado de Fuerza Capilar
    Geográfica ($FC_H$).}
    \label{fig:barplot-macro-hubs}
    \fuentefigura{Elaboración propia con \texttt{VFTModel} a partir de los resultados analíticos del modelo.}
\end{figure}

Por último, la distribución espacial y la cobertura geográfica de estos nodos de
transferencia masiva dentro del entorno urbano de la \zmvm{} se presenta mapeada
de forma detallada en la Figura~\ref{fig:mapa-fuerza-capilar}.

\begin{figure}[H]
    \centering
    \includegraphics[width=0.85\textwidth]{Figures/Cap5/FuerzaCapilar/Mapa/NB03_mapa01_macrohubs.png}
    \caption[Mapa de Macro-Hubs y Fuerza Capilar Geográfica]{
    Representación geoespacial de la \termino{Fuerza Capilar Geográfica}
    ($FC_H$). El mapa destaca la concentración de flujos de transporte en los
    principales nodos y zonas de transferencia física (Macro-Hubs) de la \zmvm.}
    \label{fig:mapa-fuerza-capilar}
    \fuentefigura{Elaboración propia con \texttt{VFTModel} utilizando visualización geoespacial.}
\end{figure}

% =============================================================================
%  5.2.1.4 FUERZA CAPILAR SISTÉMICA (MATEMÁTICA)
% =============================================================================
\subsection{\bajoRevision{Fuerza Capilar Sistémica (Matemática)}}
\label{subsec:fuerza-capilar-matematica}

\textbf{B. Fuerza Capilar Sistémica o Matemática ($FC$):} A diferencia de la
aproximación geográfica orientada a macro-hubs, esta perspectiva evalúa de manera
estricta la convergencia de flujos directamente sobre cada nodo topológico
individual. El análisis permite identificar los puntos exactos de la
infraestructura matemática de la red donde se concentra la mayor densidad de
aristas de entrada y salida, aislando singularidades críticas en la traza sin
aplicar tolerancias de proximidad espacial ni fusión de estaciones colindantes.

El Cuadro~\ref{tab:fuerza-capilar-matematica} detalla el comportamiento de los
cinco nodos individuales que registraron la mayor Fuerza Capilar pura dentro del
grafo de transporte metropolitano.

\begin{table}[H]
    \centering
    \caption[Top 5 Nodos con mayor Fuerza Capilar Sistémica (Matemática)]{
    Los cinco nodos individuales con mayor Fuerza Capilar Sistémica (Matemática)
    en la red metropolitana, evaluados de manera estricta sin agrupación
    geográfica. Corrida: 2026-04-18.}
    \label{tab:fuerza-capilar-matematica}
    \fontsize{9.5pt}{11.5pt}\selectfont
    \begin{tabular}{|c|l|c|c|c|}
        \hline
        \textbf{ID Nodo} & \textbf{Estación (Nodo Topológico)} & \textbf{Sistemas} & \textbf{Conexiones} \textit{(In / Out)} & \textbf{$FC$} \\
        \hline
        7270  & Luis Murillo                          & RTP & 21 / 20 & \textbf{41} \\
        \hline
        2817  & Calz. de Tlalpan -- Glorieta Huipulco & RTP & 19 / 15 & \textbf{34} \\
        \hline
        10211 & Sobre Calz. de Tlalpan y Luis Murillo & CC  & 16 / 17 & \textbf{33} \\
        \hline
        10717 & Tlalpan y Luis Murillo                & CC  & 18 / 15 & \textbf{33} \\
        \hline
        5540  & Estadio Azteca 2                      & RTP & 14 / 18 & \textbf{32} \\
        \hline
    \end{tabular}
    \vspace{1ex}
    \fuentefigura{Elaboración propia con \texttt{VFTModel} y datos de \texttt{Apimetro}.}
\end{table}

Las magnitudes y diferencias relativas entre los diez elementos individuales con
mayor concentración de conexiones se visualizan en el gráfico de barras de la
Figura~\ref{fig:barplot-nodos-matematicos}.

\begin{figure}[H]
    \centering
    \includegraphics[width=0.85\textwidth]{Figures/Cap5/FuerzaCapilar/Tabla/NB03_fig02_top10_nodos_matematicos_barplot.png}
    \caption[Distribución de los principales nodos individuales por Fuerza Capilar]{
    Gráfico de barras comparativo de los diez nodos individuales con mayor
    Fuerza Capilar Sistémica (Matemática) en la red metropolitana.}
    \label{fig:barplot-nodos-matematicos}
    \fuentefigura{Elaboración propia con \texttt{VFTModel} a partir de los resultados analíticos del modelo.}
\end{figure}

La distribución geográfica e implicación espacial de estos puntos específicos de
alta carga conectiva dentro de la red urbana consolidada se presenta en la
Figura~\ref{fig:mapa-nodos-matematicos}.

\begin{figure}[H]
    \centering
    \includegraphics[width=0.85\textwidth]{Figures/Cap5/FuerzaCapilar/Mapa/NB03_mapa02_nodos_matematicos.png}
    \caption[Mapa de Nodos Individuales y Fuerza Capilar Sistémica]{
    Representación geoespacial de la Fuerza Capilar Sistémica (Matemática) a
    nivel de nodo individual en la \zmvm{}.}
    \label{fig:mapa-nodos-matematicos}
    \fuentefigura{Elaboración propia con \texttt{VFTModel} utilizando visualización geoespacial.}
\end{figure}

\subsubsection{Análisis del Eje Periférico — Infraestructura Preexistente}
\label{subsubsec:periferico-fc}

Como diagnóstico previo a la propuesta del corredor anillar, se aislaron los nodos
y Macro-Hubs cuyos nombres contienen la denominación ``Periférico'', con el objetivo
de identificar los puntos de mayor tensión vehicular e intermodal sobre esta
vialidad fundamental. Este análisis se divide en dos enfoques: la concentración
geográfica (transbordos físicos en Macro-Hubs) y la concentración matemática
estricta (nodos topológicos individuales).

\paragraph{A. Fuerza Capilar Geográfica (Macro-Hubs) en el Periférico}

Al aplicar un radio de tolerancia de 100~m para agrupar las estaciones colindantes
sobre el eje, emergen los principales complejos de transbordo y alimentación. El
Cuadro~\ref{tab:periferico-hubs-geo} detalla los cinco Macro-Hubs de mayor Fuerza
Capilar Geográfica ($FC_H$), destacando la importancia de Barranca del Muerto y la
intersección con la zona de hospitales/Tlalpan como los anclajes más fuertes del
corredor.

\begin{table}[H]
    \centering
    \caption[Top 5 Macro-Hubs con mayor Fuerza Capilar sobre el Eje Periférico]{
    Cinco Macro-Hubs de transferencia con mayor Fuerza Capilar Geográfica
    ($FC_H$) ubicados sobre el Anillo Periférico. Corrida: 2026-04-18.}
    \label{tab:periferico-hubs-geo}
    \fontsize{8.5pt}{10.5pt}\selectfont
    \begin{tabularx}{\textwidth}{|c|X|c|c|c|c|}
        \hline
        \textbf{ID Hub} & \textbf{Principales Estaciones Agrupadas} & \textbf{Sistemas} & \textbf{Nodos} & \textbf{Conexiones} \textit{(In / Out)} & \textbf{$FC_H$} \\
        \hline
        259 & Nodo Barranca del Muerto (Barranca del Muerto, Alpes, Periférico Sur, Miguel Ocaranza) & CC, RTP & 21 & 50 / 42 & \textbf{92} \\
        \hline
        374 & Intersección Tlalpan--Periférico (Renato Leduc, Huipulco, Periférico, Av. Pedregal) & CC, RTP & 18 & 34 / 33 & \textbf{67} \\
        \hline
        453 & Nodo San Antonio (San Antonio, Periférico Sur, Puente Periférico Norte) & CC, RTP & 7 & 23 / 23 & \textbf{46} \\
        \hline
        402 & Eje Oriente Tepalcates (Constitución de Apatzingán, Canal de San Juan, Tepalcates) & CC, MB, RTP & 8 & 20 / 26 & \textbf{46} \\
        \hline
        291 & Intersección Xochimilco--Tepepan (16 de Septiembre, Kioto, Tec de Monterrey, Periférico Participación) & CC, RTP, TL & 14 & 23 / 22 & \textbf{45} \\
        \hline
    \end{tabularx}
    \vspace{1ex}
    \fuentefigura{Elaboración propia con \texttt{VFTModel} y datos de \texttt{Apimetro}.}
\end{table}

La Figura~\ref{fig:barplot-hubs-periferico} compara las magnitudes relativas de
estos agrupamientos geográficos, mientras que la Figura~\ref{fig:mapa-hubs-periferico}
espacializa su huella sobre la vialidad, revelando los polos naturales de atracción
de demanda que la futura infraestructura anillar deberá integrar.

\begin{figure}[H]
    \centering
    \includegraphics[width=0.85\textwidth]{Figures/Cap5/FuerzaCapilar/Tabla/NB03_fig04_top10_macrohubs_periferico_barplot.png}
    \caption[Fuerza Capilar de Macro-Hubs en Anillo Periférico]{Gráfico de barras
    de los diez principales Macro-Hubs geográficos ubicados en el Anillo Periférico,
    ordenados por su nivel de Fuerza Capilar.}
    \label{fig:barplot-hubs-periferico}
    \fuentefigura{Elaboración propia con \texttt{VFTModel}.}
\end{figure}

\begin{figure}[H]
    \centering
    \includegraphics[width=0.85\textwidth]{Figures/Cap5/FuerzaCapilar/Mapa/NB03_mapa03_macrohubs_periferico.png}
    \caption[Distribución espacial de Macro-Hubs en el Periférico]{Mapa de la
    Fuerza Capilar Geográfica ($FC_H$) en los Macro-Hubs del Anillo Periférico.}
    \label{fig:mapa-hubs-periferico}
    \fuentefigura{Elaboración propia con \texttt{VFTModel} utilizando visualización geoespacial.}
\end{figure}

\paragraph{B. Fuerza Capilar Sistémica (Nodos Matemáticos) en el Periférico}

Al desagregar los Hubs y analizar la Fuerza Capilar Sistémica ($FC$) estricta a
nivel de nodo individual, es posible identificar los paraderos específicos que
soportan la mayor carga de conexiones en la infraestructura actual. El
Cuadro~\ref{tab:periferico-nodos-math} presenta los cinco paraderos matemáticos
de mayor convergencia de flujo.

\begin{table}[H]
    \centering
    \caption[Nodos con mayor Fuerza Capilar Sistémica sobre el Eje Periférico]{
    Cinco nodos topológicos individuales con mayor Fuerza Capilar Matemática
    ($FC$) ubicados sobre el Anillo Periférico. Corrida: 2026-04-18.}
    \label{tab:periferico-nodos-math}
    \fontsize{9.5pt}{11.5pt}\selectfont
    \begin{tabular}{|c|l|c|r|r|c|}
        \hline
        \textbf{ID Nodo} & \textbf{Estación (Nodo Topológico)} & \textbf{Sistema} & \textbf{Entr.} & \textbf{Sal.} & \textbf{$FC$} \\
        \hline
        707  & Anillo Periférico y Blvd. Adolfo Ruiz Cortines & RTP & 14 & 16 & \textbf{30} \\
        \hline
        706  & Anillo Periférico y Blvd. Adolfo Ruiz Cortines & RTP & 13 & 13 & \textbf{26} \\
        \hline
        8538 & Periférico -- Barranca del Muerto              &  CC & 12 &  9 & \textbf{21} \\
        \hline
        602  & Anillo Periférico Canal de Garay -- Bilbao     & RTP & 10 & 10 & \textbf{20} \\
        \hline
        8510 & Periférico                                     & RTP &  9 & 11 & \textbf{20} \\
        \hline
    \end{tabular}
    \vspace{1ex}
    \fuentefigura{Elaboración propia con \texttt{VFTModel} y datos de \texttt{Apimetro}.}
\end{table}

La jerarquía analítica de estos puntos singulares se constata en la
Figura~\ref{fig:barplot-nodos-periferico}, seguida de la
Figura~\ref{fig:mapa-nodos-periferico}, que detalla la distribución geográfica a
escala de paradero individual, útil para el micro-posicionamiento de futuras
estaciones anillares.

\begin{figure}[H]
    \centering
    \includegraphics[width=0.85\textwidth]{Figures/Cap5/FuerzaCapilar/Tabla/NB03_fig05_top10_nodos_matematicos_periferico_barplot.png}
    \caption[Fuerza Capilar de nodos individuales en Anillo Periférico]{
    Gráfico de barras comparativo de los diez principales nodos matemáticos
    individuales ubicados en el Anillo Periférico.}
    \label{fig:barplot-nodos-periferico}
    \fuentefigura{Elaboración propia con \texttt{VFTModel}.}
\end{figure}

\begin{figure}[H]
    \centering
    \includegraphics[width=0.85\textwidth]{Figures/Cap5/FuerzaCapilar/Mapa/NB03_mapa04_nodos_matematicos_periferico.png}
    \caption[Distribución espacial de nodos matemáticos en el Periférico]{Mapa
    de la Fuerza Capilar Sistémica (Matemática) a nivel de nodo individual sobre
    el corredor del Anillo Periférico.}
    \label{fig:mapa-nodos-periferico}
    \fuentefigura{Elaboración propia con \texttt{VFTModel} utilizando visualización geoespacial.}
\end{figure}

\vspace{2ex}
\begin{quote}
    \small\textit{Nota aclaratoria:} los datos presentados en esta sección
    reflejan exclusivamente la infraestructura de transporte \textbf{preexistente}
    sobre el Anillo Periférico. No se ha trazado ni modelado aún la geometría
    del corredor anillar propuesto; la sección permanece en plano especulativo
    para identificar los puntos naturales de anclaje de la futura línea a nivel
    capilar.
\end{quote}

% !TeX root = ../../tesis.tex
% =============================================================================
%  5.2.1.5 FACTOR DE DESVIACIÓN
% =============================================================================
\subsection{\bajoRevision{Factor de Desviación ($DF$)}}
\label{subsec:detour-factor}
\notaRevision{Los resultados de esta subsección están sujetos a revisión con el
comité tutor. Los valores presentados corresponden a una muestra de 100~rutas
aleatorias ejecutadas el 27 de abril de 2026.}

El \textbf{Factor de Desviación} ($DF$) mide la eficiencia topológica de la
red comparando la distancia real recorrida a través de la infraestructura
contra la distancia teórica en línea recta (Haversine) entre el mismo par
origen-destino:

\begin{equation}
    DF = \frac{D_{\text{red}} + D_{\text{caminata}}}{D_{\text{haversine}}(O, D)}
    \label{eq:detour-factor-res}
\end{equation}

\begin{itemize}
    \item $D_{\text{red}}$: distancia acumulada sobre el trazo real de la ruta
    óptima en tiempo.
    \item $D_{\text{caminata}}$: suma de la primera y última milla peatonal.
    \item $D_{\text{haversine}}(O,D)$: distancia en línea recta entre origen y
    destino calculada con la fórmula de Haversine.
\end{itemize}

Un valor de $DF = 1.0$ indica una ruta perfectamente directa; valores superiores
a $1.5$ revelan ineficiencias topológicas significativas.

% -----------------------------------------------------------------------------
\subsubsection{Estadísticas Globales de la Red}
\label{subsubsec:df-estadisticas}

Sobre una muestra aleatoria de \textbf{100 pares origen-destino} (semilla
reproducible $= 42$), el análisis arrojó los siguientes estadísticos globales:

\begin{table}[H]
    \centering
    \caption[Estadísticos globales del Factor de Desviación — muestra de 100 rutas]{
    Estadísticos descriptivos del $DF$ sobre 100 pares O-D de la red de la \zmvm.
    Corrida: 2026-04-27.}
    \label{tab:df-estadisticas}
    \begin{tabular}{|l|r|}
        \hline
        \textbf{Estadístico} & \textbf{Valor} \\
        \hline
        $n$ (rutas calculadas)     & 100   \\
        Media ($\bar{DF}$)         & 1.545 \\
        Desviación estándar        & 0.517 \\
        Mínimo                     & 0.97  \\
        Percentil 25               & 1.27  \\
        Mediana                    & 1.40  \\
        Percentil 75               & 1.67  \\
        Máximo                     & 4.34  \\
        \hline
    \end{tabular}
    \fuentefigura{Elaboración propia con \texttt{VFTModel} (corrida 2026-04-27).}
\end{table}

La Figura~\ref{fig:df-caso-inicial} muestra la distribución geoespacial de los
100 pares O-D y las rutas calculadas sobre el territorio de la \zmvm{}.

\begin{figure}[H]
    \centering
    \includegraphics[width=0.85\textwidth]{Figures/Cap5/FactorDesviacion/Mapa/NB04_mapa01_caso_inicial.png}
    \caption[Distribución de la muestra O-D para el Factor de Desviación]{
    Distribución de los 100 pares origen-destino de la muestra aleatoria y sus
    rutas óptimas en la \zmvm. Corrida: 2026-04-27.}
    \label{fig:df-caso-inicial}
    \fuentefigura{Elaboración propia con \texttt{VFTModel} (corrida 2026-04-27).}
\end{figure}

% -----------------------------------------------------------------------------
\subsubsection{Casos de Estudio Seleccionados}
\label{subsubsec:df-casos}

Se analizaron doce rutas que evalúan la conectividad desde las periferias sur y
oriente hacia distintos destinos de la metrópoli. Las
Figuras~\ref{fig:df-casos-1-4}--\ref{fig:df-casos-9-12} presentan los trazados
cartográficos de cada caso; el Cuadro~\ref{tab:df-casos} consolida los valores
numéricos al final de esta sección.

% --- Casos 1–4 ---------------------------------------------------------------
\begin{figure}[H]
    \centering
    \begin{subfigure}[t]{0.48\textwidth}
        \includegraphics[width=\textwidth]{Figures/Cap5/FactorDesviacion/Mapa/NB04_mapa04_caso01.png}
        \caption{Caso~1: Milpa Alta → Autódromo. $DF = 1.29$.}
        \label{fig:df-c01}
    \end{subfigure}
    \hfill
    \begin{subfigure}[t]{0.48\textwidth}
        \includegraphics[width=\textwidth]{Figures/Cap5/FactorDesviacion/Mapa/NB04_mapa05_caso02.png}
        \caption{Caso~2: Milpa Alta → Reforma (CBD). $DF = 1.37$.}
        \label{fig:df-c02}
    \end{subfigure}

    \vspace{0.5em}
    \begin{subfigure}[t]{0.48\textwidth}
        \includegraphics[width=\textwidth]{Figures/Cap5/FactorDesviacion/Mapa/NB04_mapa06_caso03.png}
        \caption{Caso~3: Milpa Alta → Santa Fe (CBD). $DF = 1.27$.}
        \label{fig:df-c03}
    \end{subfigure}
    \hfill
    \begin{subfigure}[t]{0.48\textwidth}
        \includegraphics[width=\textwidth]{Figures/Cap5/FactorDesviacion/Mapa/NB04_mapa07_caso04.png}
        \caption{Caso~4: Milpa Alta → FES Acatlán. $DF = 1.35$.}
        \label{fig:df-c04}
    \end{subfigure}
    \caption[Casos de estudio del Factor de Desviación: 1–4]{Rutas de los casos 1 al 4.
    Origen común: Milpa Alta hacia distintos destinos del norte y poniente.}
    \label{fig:df-casos-1-4}
    \fuentefigura{Elaboración propia con \texttt{VFTModel} (corrida 2026-04-27).}
\end{figure}

% --- Casos 5–8 ---------------------------------------------------------------
\begin{figure}[H]
    \centering
    \begin{subfigure}[t]{0.48\textwidth}
        \includegraphics[width=\textwidth]{Figures/Cap5/FactorDesviacion/Mapa/NB04_mapa08_caso05.png}
        \caption{Caso~5: Milpa Alta → C.U. $DF = 1.32$.}
        \label{fig:df-c05}
    \end{subfigure}
    \hfill
    \begin{subfigure}[t]{0.48\textwidth}
        \includegraphics[width=\textwidth]{Figures/Cap5/FactorDesviacion/Mapa/NB04_mapa09_caso06.png}
        \caption{Caso~6: Milpa Alta → Coyoacán. $DF = 1.34$.}
        \label{fig:df-c06}
    \end{subfigure}

    \vspace{0.5em}
    \begin{subfigure}[t]{0.48\textwidth}
        \includegraphics[width=\textwidth]{Figures/Cap5/FactorDesviacion/Mapa/NB04_mapa10_caso07.png}
        \caption{Caso~7: Santa Catarina → C.U. $DF = 1.82$.}
        \label{fig:df-c07}
    \end{subfigure}
    \hfill
    \begin{subfigure}[t]{0.48\textwidth}
        \includegraphics[width=\textwidth]{Figures/Cap5/FactorDesviacion/Mapa/NB04_mapa11_caso08.png}
        \caption{Caso~8: Santa Catarina → FES Acatlán. $DF = 1.42$.}
        \label{fig:df-c08}
    \end{subfigure}
    \caption[Casos de estudio del Factor de Desviación: 5–8]{Rutas de los casos 5 al 8.
    Destinos en el sur de la ciudad y primeros casos desde Santa Catarina.}
    \label{fig:df-casos-5-8}
    \fuentefigura{Elaboración propia con \texttt{VFTModel} (corrida 2026-04-27).}
\end{figure}

% --- Casos 9–12 --------------------------------------------------------------
\begin{figure}[H]
    \centering
    \begin{subfigure}[t]{0.48\textwidth}
        \includegraphics[width=\textwidth]{Figures/Cap5/FactorDesviacion/Mapa/NB04_mapa12_caso09.png}
        \caption{Caso~9: Milpa Alta → Santa Catarina. $DF = 2.52$.}
        \label{fig:df-c09}
    \end{subfigure}
    \hfill
    \begin{subfigure}[t]{0.48\textwidth}
        \includegraphics[width=\textwidth]{Figures/Cap5/FactorDesviacion/Mapa/NB04_mapa13_caso10.png}
        \caption{Caso~10: Milpa Alta → Ajusco. $DF = 1.42$.}
        \label{fig:df-c10}
    \end{subfigure}

    \vspace{0.5em}
    \begin{subfigure}[t]{0.48\textwidth}
        \includegraphics[width=\textwidth]{Figures/Cap5/FactorDesviacion/Mapa/NB04_mapa14_caso11.png}
        \caption{Caso~11: Ajusco → Preparatoria 1. $DF = 1.43$.}
        \label{fig:df-c11}
    \end{subfigure}
    \hfill
    \begin{subfigure}[t]{0.48\textwidth}
        \includegraphics[width=\textwidth]{Figures/Cap5/FactorDesviacion/Mapa/NB04_mapa15_caso12.png}
        \caption{Caso~12: P. Los Olivos → Centro Tláhuac. $DF = 1.06$.}
        \label{fig:df-c12}
    \end{subfigure}
    \caption[Casos de estudio del Factor de Desviación: 9–12]{Rutas de los casos 9 al 12.
    El caso~9 (sur-oriente transversal) registra el mayor $DF$ del conjunto.}
    \label{fig:df-casos-9-12}
    \fuentefigura{Elaboración propia con \texttt{VFTModel} (corrida 2026-04-27).}
\end{figure}

% --- Extremos globales de la muestra aleatoria -------------------------------
Para situar los casos de estudio en el contexto de la muestra completa, las
Figuras~\ref{fig:df-mejor-caso} y~\ref{fig:df-peor-caso} ilustran los extremos
globales obtenidos en los 100 pares O-D aleatorios.

\begin{figure}[H]
    \centering
    \begin{subfigure}[t]{0.48\textwidth}
        \includegraphics[width=\textwidth]{Figures/Cap5/FactorDesviacion/Mapa/NB04_mapa02_ruta_mas_eficiente.png}
        \caption{Mejor caso: Plutarco Elías Calles → Fuente de Loreto. $DF = 0.97$.}
        \label{fig:df-mejor-caso}
    \end{subfigure}
    \hfill
    \begin{subfigure}[t]{0.48\textwidth}
        \includegraphics[width=\textwidth]{Figures/Cap5/FactorDesviacion/Mapa/NB04_mapa03_ruta_mas_ineficiente.png}
        \caption{Peor caso: Carlos A. Vidal → 20 de Noviembre. $DF = 4.34$.}
        \label{fig:df-peor-caso}
    \end{subfigure}
    \caption[Extremos del Factor de Desviación en la muestra aleatoria]{
    Rutas con menor y mayor Factor de Desviación en la muestra aleatoria de 100 pares O-D.}
    \label{fig:df-extremos}
    \fuentefigura{Elaboración propia con \texttt{VFTModel} (corrida 2026-04-27).}
\end{figure}

% --- Tabla resumen de los 12 casos -------------------------------------------
\begin{table}[H]
    \centering
    \caption[Factor de Desviación — resumen de los doce casos de estudio]{
    Resultados del Factor de Desviación para los doce casos estratégicos de la
    \zmvm{} ilustrados en las figuras anteriores. Corrida: 2026-04-27.}
    \label{tab:df-casos}
    \small
    \begin{tabularx}{\textwidth}{|l|X|X|r|r|r|}
        \hline
        \textbf{\#} & \textbf{Origen} & \textbf{Destino} &
        \textbf{$D_{\text{red}}$ (km)} & \textbf{$D_{\text{Hav.}}$ (km)} & \textbf{$DF$} \\
        \hline
         1 & Milpa Alta        & Autódromo / Centro-Oriente   & 29.18 & 22.59 & 1.29 \\
         2 & Milpa Alta        & Paseo de la Reforma (CBD)    & 37.73 & 27.50 & 1.37 \\
         3 & Milpa Alta        & Santa Fe (CBD)               & 36.42 & 28.62 & 1.27 \\
         4 & Milpa Alta        & FES Acatlán                  & 48.34 & 35.73 & 1.35 \\
         5 & Milpa Alta        & Ciudad Universitaria         & 26.25 & 19.96 & 1.32 \\
         6 & Milpa Alta        & Coyoacán                     & 25.24 & 18.86 & 1.34 \\
         7 & Santa Catarina    & Ciudad Universitaria         & 42.06 & 23.10 & 1.82 \\
         8 & Santa Catarina    & FES Acatlán                  & 45.64 & 32.16 & 1.42 \\
         9 & Milpa Alta        & Santa Catarina (Sur-Oriente) & 37.22 & 14.79 & 2.52 \\
        10 & Milpa Alta        & Ajusco                       & 29.19 & 20.56 & 1.42 \\
        11 & Ajusco            & Preparatoria 1               & 15.91 & 11.16 & 1.43 \\
        12 & Parque Los Olivos & Centro de Tláhuac            &  2.20 &  2.08 & 1.06 \\
        \hline
        \multicolumn{5}{|l|}{\textbf{Mejor caso global (muestra aleatoria):}
        Plutarco Elias Calles → Fuente de Loreto} & \textbf{0.97} \\
        \multicolumn{5}{|l|}{\textbf{Peor caso global (muestra aleatoria):}
        Carlos A. Vidal → 20 de Noviembre} & \textbf{4.34} \\
        \hline
    \end{tabularx}
    \fuentefigura{Elaboración propia con \texttt{VFTModel} (corrida 2026-04-27).}
\end{table}

El análisis revela que los viajes periféricos sur-oriente (Caso~9: $DF = 2.52$,
Figura~\ref{fig:df-c09}) registran la mayor ineficiencia entre los casos de
estudio, obligando al usuario a recorrer 37.2~km de red para cubrir una distancia
lineal de 14.8~km. Esta situación responde a la ausencia de conectividad
transversal en el arco sur-oriente de la ciudad, que impone al tránsito
perimetral el ingreso al centro metropolitano como nodo obligado, confirmando la
hipótesis sobre la necesidad de un corredor anillar en el Periférico.


% !TeX root = ../tesis.tex
%%%%%%%%%%%%%%%%%%%%%%%%%%%%%%%%%%%%%%%%%%%%%%%%%%%%%%%%%%%%%%%%%%%
% SECCIÓN 5.2-2 — RESULTADOS VFT: FASE 2
% Ponderación dinámica de la red
%%%%%%%%%%%%%%%%%%%%%%%%%%%%%%%%%%%%%%%%%%%%%%%%%%%%%%%%%%%%%%%%%%%
\section{\bajoRevision{Fase 2: Ponderación Dinámica de la Red}}
\label{sec:resultados-fase2}

\notaRevision{Los indicadores de esta fase se encuentran en revisión con el
comité tutor. Los coeficientes de fricción vial y la penalización por
transferencia han sido calculados e integrados al grafo; el \termino{Node Score}
se encuentra en etapa de planificación, pendiente de la disponibilidad de las
vistas materializadas MV1 y MV2 de \textbf{Apimetro}.}

La Fase~2 del Modelo VFT enriquece el grafo construido en la Fase~1 con dos
capas de pesos dinámicos que aproximan las condiciones reales de operación de la
red: el \textbf{Coeficiente de Fricción Vial} ($CF$), que penaliza el tiempo de
traslado en función del nivel de interferencia del tráfico urbano sobre cada
tipo de vía, y la \textbf{Penalización por Transferencia} ($T_b$), que incorpora
el costo de caminata y espera en cada transbordo intermodal. Estos dos pesos
conforman el componente dinámico de la \textbf{Ecuación VFT}:

\begin{equation}
    w(u,v) = \frac{D_{uv}}{V_{\text{modo}}} \cdot CF_{\text{vía}} + T_b
    \label{eq:vft-arista}
\end{equation}

\begin{itemize}
    \item $D_{uv}$: distancia geodésica del segmento calculada mediante la
    fórmula de \termino{Haversine} (metros).
    \item $V_{\text{modo}}$: velocidad comercial ajustada del sistema de
    transporte (km/h), extraída del \geojson{} de \textbf{Apimetro} o
    sustituida por el valor de \termino{fallback} calibrado por sistema.
    \item $CF_{\text{vía}}$: coeficiente de fricción vial según el tipo de
    derecho de vía del segmento (adimensional, $\geq 1$).
    \item $T_b$: penalización por transferencia peatonal en minutos; se aplica
    exclusivamente en las aristas de transbordo generadas por el
    \termino{Snapping Lógico Geográfico}.
\end{itemize}

Para las aristas de servicio normal, el peso de enrutamiento es
$\mathit{weight} = (D_{uv}/V_{\text{modo}}) \cdot CF_{\text{vía}}$; el
$T_b$ queda registrado como atributo informativo pero no penaliza el
segmento. La penalización por transferencia actúa únicamente sobre aristas de
tipo \texttt{transfer}, como se describe en la Sección~\ref{subsec:res-penalizacion-transferencia}.

% ----------------------------------------------------------------
\subsection{Coeficiente de Fricción Vial ($CF$)}
\label{subsec:coeficiente-friccion}
% ----------------------------------------------------------------

El \textbf{Coeficiente de Fricción Vial} ($CF$) traduce el tipo de
\termino{derecho de vía} de cada segmento en un penalizador multiplicativo
sobre el tiempo de traslado. La premisa es que circular por un túnel del
Metro —aislado del tráfico en superficie— implica un costo temporal muy
distinto al de un camión de la \textsc{rtp} inmerso en el flujo vehicular de
Insurgentes o Periférico.

\subsubsection{Constante de Saturación Vial}
\label{subsubsec:beta-saturacion}

El parámetro base del coeficiente es la constante de saturación $\beta$,
derivada del \termino{TomTom Traffic Index} para la \zmvm:

\begin{equation}
    \beta_{\text{ZMVM}} = 0.759
    \label{eq:beta-zmvm}
\end{equation}

Este valor indica que la ciudad opera, en condiciones promedio, al 76~\% de
su capacidad vial. Es decir, un vehículo en tráfico mixto experimenta, en
promedio, un sobretiempo del 76~\% respecto al recorrido en flujo libre.

\subsubsection{Fórmula y Valores por Derecho de Vía}
\label{subsubsec:formula-cf}

El coeficiente $CF$ se obtiene a partir de un parámetro de escala $\alpha$
que modula la exposición de cada tipo de vía a la saturación urbana:

\begin{equation}
    CF = 1 + \alpha \cdot \beta_{\text{ZMVM}}
    \label{eq:cf-formula}
\end{equation}

\begin{itemize}
    \item $\alpha \in \{0.0,\ 0.2,\ 0.5,\ 1.0\}$: fracción de la saturación
    urbana que afecta el segmento, determinada por su tipo de derecho de vía.
    \item $\beta_{\text{ZMVM}} = 0.759$: constante de saturación de la \zmvm
    (\termino{TomTom Traffic Index}).
    \item $CF \geq 1$: un valor de 1.0 indica flujo ideal sin interferencia;
    valores superiores representan el penalizador proporcional al congestionamiento.
\end{itemize}

El Cuadro~\ref{tab:cf-por-via} sintetiza los cuatro tipos de derecho de vía
modelados, los sistemas de transporte correspondientes y el valor numérico del
coeficiente resultante.

\begin{table}[H]
    \centering
    \caption[Coeficiente de Fricción Vial por tipo de derecho de vía]{Coeficiente
    de Fricción Vial ($CF$) por tipo de derecho de vía en el Modelo VFT.
    $\alpha$ escala el impacto de la saturación urbana
    $\beta_{\text{ZMVM}} = 0.759$.}
    \label{tab:cf-por-via}
    \begin{tabularx}{\textwidth}{|l|c|c|X|}
        \hline
        \textbf{Derecho de vía} & \textbf{$\alpha$} & \textbf{$CF$} &
        \textbf{Sistemas representativos} \\
        \hline
        Exclusivo   & 0.0 & 1.000 & Metro, Tren Suburbano \\
        Confinado   & 0.2 & 1.152 & Metrobús, Cablebús, Mexicable \\
        Compartido  & 0.5 & 1.380 & Trolebús, Tren Ligero \\
        Mixto       & 1.0 & 1.759 & \textsc{rtp}, Corredores Concesionados,
                                    Mexibús, Pumabús \\
        \hline
    \end{tabularx}
    \fuentefigura{Elaboración propia con base en \texttt{impedance.py}
    (VFTModel, corrida 2026-04-18). Los valores $\alpha$ tienen como referencia
    normativa el \textit{Manual de Calles: Diseño Vial para Ciudades Mexicanas}
    \parencite{sedatu2019manual}.}
\end{table}

El rango del coeficiente $CF \in [1.000,\ 1.759]$ implica que el tiempo de
traslado en un segmento de vía mixta (RTP, Corredores Concesionados) puede ser
hasta un 75.9~\% mayor que el mismo recorrido bajo condiciones de flujo ideal.
Esta diferencia no refleja únicamente la velocidad operacional del vehículo —que
ya está calibrada en $V_{\text{modo}}$— sino el efecto sistemático de las
detenciones por semáforos, cruces peatonales y congestión que son inherentes a
los segmentos de superficie compartida.

En el caso del sistema de derecho de vía desconocido, el Modelo VFT asigna
por defecto el valor de vía \texttt{mixto} ($CF = 1.759$), aplicando el
principio de escenario más conservador para no subestimar los tiempos de traslado.

% ----------------------------------------------------------------
\subsection{Penalización por Transferencia ($T_b$)}
\label{subsec:res-penalizacion-transferencia}
% ----------------------------------------------------------------

La \textbf{Penalización por Transferencia} ($T_b$) modela el costo temporal que
un pasajero incurre al cambiar de sistema o línea dentro de la red. A diferencia
del tiempo de traslado —que depende de la distancia y la velocidad del vehículo—,
el costo de transbordo consta de dos componentes independientes: el tiempo de
caminata entre plataformas y el tiempo de espera en el andén de destino.

\subsubsection{Fórmula de la Penalización}
\label{subsubsec:formula-tb}

Para cada arista de transbordo generada por el \termino{Snapping Lógico Geográfico}
entre los nodos $u$ (origen del transbordo) y $v$ (destino del transbordo), el
peso de la arista es:

\begin{equation}
    T_b(u, v) = t_{\text{caminata}}(u, v) + \frac{F_v}{2}
    \label{eq:tb-formula}
\end{equation}

\begin{itemize}
    \item $t_{\text{caminata}}(u, v)$: tiempo de caminata entre los dos nodos
    a una velocidad peatonal normativa de 5~km/h (1.38~m/s), equivalente al
    límite superior establecido en el \textit{Manual de Calles SEDATU}
    \parencite{sedatu2019manual} para infraestructura de transbordo sin detenciones
    en semáforo.
    \item $F_v$: frecuencia de servicio del sistema destino $v$ (minutos entre
    unidades); si no está reportada en el \geojson{}, se utiliza el valor
    de \termino{fallback} calibrado por sistema.
    \item $F_v / 2$: tiempo de espera esperado, calculado como la mitad de la
    frecuencia bajo la hipótesis de llegada aleatoria del pasajero al andén
    (\termino{random arrival assumption}).
\end{itemize}

El tiempo de caminata se calcula como:

\begin{equation}
    t_{\text{caminata}}(u, v) = \frac{D_{uv}}{V_{\text{peatonal}}}
    = \frac{D_{uv}}{83.33\ \text{m/min}}
    \label{eq:t-caminata}
\end{equation}

donde $D_{uv}$ es la distancia geodésica entre los dos nodos (Haversine) en
metros, y $V_{\text{peatonal}} = 5000\ \text{m} / 60\ \text{min} = 83.33\ \text{m/min}$.

\subsubsection{Valores de Frecuencia de Fallback}
\label{subsubsec:fallback-frecuencia}

El Cuadro~\ref{tab:fallback-frecuencia} presenta los valores de frecuencia de
\termino{fallback} calibrados para cada sistema de transporte, junto con el
costo de abordaje resultante ($F_v / 2$). Estos valores representan intervalos
promedio en condiciones de operación normal en hora no pico.

\begin{table}[H]
    \centering
    \caption[Frecuencias de \textit{fallback} y costo de abordaje por sistema]{Frecuencias
    de servicio de \termino{fallback} y penalización de abordaje resultante
    ($T_b^{\text{espera}} = F_v / 2$) por sistema de transporte en el Modelo VFT.}
    \label{tab:fallback-frecuencia}
    \begin{tabular}{|l|c|c|}
        \hline
        \textbf{Sistema} & \textbf{$F_v$ (min)} & \textbf{$F_v / 2$ (min)} \\
        \hline
        \textsc{Metro} (STC)          & 3.0  & 1.50 \\
        Metrobús (MB)                 & 5.0  & 2.50 \\
        Tren Ligero (TL)              & 10.0 & 5.00 \\
        Trolebús (TROLE)              & 8.0  & 4.00 \\
        \textsc{rtp}                  & 15.0 & 7.50 \\
        Tren Suburbano (SUB)          & 12.0 & 6.00 \\
        Cablebús (CBB)                & 10.0 & 5.00 \\
        Corredores Concesionados (CC) & 15.0 & 7.50 \\
        Mexibús                       & 6.0  & 3.00 \\
        \hline
    \end{tabular}
    \fuentefigura{Elaboración propia con base en \texttt{graph\_builder.py}
    (VFTModel, corrida 2026-04-18). Los valores corresponden a frecuencias
    promedio en horario valle; las frecuencias en hora pico son reportadas por
    los operadores como atributo del \geojson{} y tienen precedencia sobre los
    \termino{fallbacks}.}
\end{table}

\subsubsection{Decisión de Diseño: Boarding Cost Informativo}
\label{subsubsec:boarding-cost-diseno}

Una precisión metodológica relevante: el costo de abordaje $F_v / 2$ \textbf{no
se suma al peso de las aristas de servicio normal}. Solo actúa como componente
del peso de las aristas de transbordo peatonal. Esta decisión responde a que el
pasajero ya está dentro del vehículo en los segmentos normales; el tiempo de
espera por frecuencia solo es pertinente en el momento de abordar un nuevo
sistema. Sumarlo en cada arista normal equivaldría a penalizar al pasajero en
cada estación intermedia, sobreestimando sistemáticamente los tiempos de ruta.

En consecuencia, el Modelo VFT registra $F_v / 2$ como atributo informativo
en las aristas de servicio (\texttt{boarding\_cost\_min}) para análisis
descriptivos, pero el atributo \texttt{weight} de estas aristas solo contiene
el tiempo de traslado con fricción.

% ----------------------------------------------------------------
\subsection{\bajoRevision{Node Score — Estado de Implementación}}
\label{subsec:node-score}
% ----------------------------------------------------------------

\notaRevision{El \termino{Node Score} ($\text{VFT\_score}$) se encuentra en etapa
de planificación. Su cálculo requiere los datos de afluencia por línea (MV1) y
los polígonos AGEB con atributos censales (MV2) de \textbf{Apimetro}, aún no
disponibles en la fecha de cierre de este documento. La arquitectura del indicador,
sus fórmulas y las dependencias de datos se presentan a continuación como
fundamento para la etapa de implementación.}

El \termino{Node Score} ($\text{VFT\_score}_i$) constituye la segunda capa de
pesos del Modelo VFT: mientras que los pesos de arista modelan el costo temporal
de recorrer la red, el \termino{Node Score} asigna a cada nodo una medida de
\textbf{importancia compuesta} que combina propiedades topológicas de la red con
información de demanda observada y potencial. Este indicador es el que alimenta
el análisis de vulnerabilidad de la Fase~3.

\subsubsection{Fórmula del Node Score}
\label{subsubsec:formula-vft-score}

\begin{equation}
    \text{VFT\_score}_i = \alpha_1 \hat{k}_i + \alpha_2 \hat{a}_i + \alpha_3 \hat{p}_i
    \label{eq:vft-score}
\end{equation}

donde los tres componentes normalizados al rango $[0, 1]$ son:

\begin{itemize}
    \item $\hat{k}_i$: \textbf{Fuerza Capilar normalizada} (min-max sobre todos
    los nodos del grafo). Calculada en la Fase~1 para los 11{,}115 nodos.
    \item $\hat{a}_i$: \textbf{Afluencia estimada por estación normalizada}.
    Derivada de la afluencia diaria por línea ($A_\ell$, fuente MV1) desagregada
    internamente por VFT mediante ponderación por grado topológico:
    \begin{equation}
        a_i = \sum_{\ell \ni i} A_\ell \cdot
        \frac{g_i}{\displaystyle\sum_{j \in \ell} g_j}
        \label{eq:desagregacion-afluencia}
    \end{equation}
    donde $g_i = \deg(i)$ es el grado del nodo $i$ en el grafo NetworkX.
    \item $\hat{p}_i$: \textbf{Población captada en buffer de 800~m normalizada}.
    Calculada intersectando un radio de 800~m alrededor de cada estación con los
    polígonos AGEB (fuente MV2), ponderando cada AGEB por la fracción de su área
    dentro del buffer:
    \begin{equation}
        p_i = \sum_{j \in \mathcal{A}_i} \text{pob}_j \cdot
        \frac{\text{área}(B_i \cap \text{AGEB}_j)}{\text{área}(\text{AGEB}_j)}
        \label{eq:poblacion-captada}
    \end{equation}
    donde $B_i$ es el buffer circular de 800~m en proyección métrica EPSG:6372
    (México Centro) y $\mathcal{A}_i$ el conjunto de AGEBs que lo intersectan.
\end{itemize}

La normalización de cada componente usa el método \termino{min-max}:
\begin{equation}
    \hat{x}_i = \frac{x_i - \min(\mathbf{x})}{\max(\mathbf{x}) - \min(\mathbf{x})}
    \label{eq:normalizacion-minmax}
\end{equation}

\subsubsection{Determinación de los Pesos $\alpha$}
\label{subsubsec:pesos-alpha}

Los pesos $\alpha_1, \alpha_2, \alpha_3$ se determinarán en dos etapas
progresivas una vez que los datos de MV1 y MV2 estén disponibles:

\begin{enumerate}
    \item \textbf{Validación inicial (pesos iguales):} $\alpha_1 = \alpha_2 =
    \alpha_3 = 1/3$. Objetivo: verificar que nodos de alta importancia reconocida
    (Pantitlán, Insurgentes, Pino Suárez, Balderas) aparecen en el top-10 del
    \termino{ranking}. Si no lo hacen, el problema se ubica en los datos de entrada.
    \item \textbf{Pesos por entropía (CRITIC):} los pesos finales se derivan de
    la variabilidad observada de cada componente sobre todos los nodos:
    \begin{equation}
        \alpha_n = \frac{\sigma_n}{\displaystyle\sum_{m=1}^{3} \sigma_m}
        \label{eq:alpha-entropia}
    \end{equation}
    donde $\sigma_n = \text{std}(\hat{x}_n)$. Un componente más discriminante
    (mayor dispersión entre nodos) recibe mayor peso, penalizando automáticamente
    los componentes redundantes.
\end{enumerate}

\subsubsection{Dependencias de Datos y Estado}
\label{subsubsec:node-score-estado}

El Cuadro~\ref{tab:node-score-dependencias} resume los componentes del
\termino{Node Score}, su fuente de datos y su estado de disponibilidad a la
fecha de cierre del documento.

\begin{table}[H]
    \centering
    \caption[Dependencias del Node Score por componente]{Componentes del
    $\text{VFT\_score}$, fuentes de datos requeridas y estado de implementación.}
    \label{tab:node-score-dependencias}
    \begin{tabularx}{\textwidth}{|l|X|l|l|}
        \hline
        \textbf{Componente} & \textbf{Fuente} & \textbf{Responsable} & \textbf{Estado} \\
        \hline
        $\hat{k}_i$ Fuerza Capilar    & VFT interno (\texttt{capillary\_strength.py})
                                        & VFTModel & Calculado \\
        $\hat{a}_i$ Afluencia         & MV1 Apimetro — afluencia por línea
                                        & Apimetro & Pendiente \\
        $\hat{p}_i$ Población captada & MV2 Apimetro — AGEBs Censo INEGI 2020
                                        & Apimetro & Pendiente \\
        Pesos $\alpha$ (entropía)     & Derivados de $\sigma(\hat{k}, \hat{a}, \hat{p})$
                                        & VFTModel & Pendiente MV1+MV2 \\
        \hline
    \end{tabularx}
    \fuentefigura{Elaboración propia con base en
    \texttt{weighted\_graph\_enrichment.md} y
    \texttt{apimetro\_requerimientos\_mvs.md} (VFTModel, 2026-04-28).}
\end{table}

El diseño del \termino{Node Score} garantiza la \textbf{independencia entre
componentes}: la desagregación de afluencia usa el grado topológico crudo
$g_i = \deg(i)$ —no la Fuerza Capilar $k_i$— para evitar que $\hat{k}_i$
domine doblemente el cómputo. Antes de asignar los pesos finales por entropía,
se verificará la correlación $\rho(\hat{k}, \hat{a})$: si $\rho \geq 0.75$,
los dos componentes se fusionarán en un índice combinado $\delta_i = \beta
\hat{k}_i + (1-\beta)\hat{a}_i$ antes de calcular el \termino{score} final.

% !TeX root = ../tesis.tex
%%%%%%%%%%%%%%%%%%%%%%%%%%%%%%%%%%%%%%%%%%%%%%%%%%%%%%%%%%%%%%%%%%%
% SECCIÓN 5.2-3 — RESULTADOS VFT: FASE 3
% Análisis de red: tiempo de viaje, centralidad y vulnerabilidad
%%%%%%%%%%%%%%%%%%%%%%%%%%%%%%%%%%%%%%%%%%%%%%%%%%%%%%%%%%%%%%%%%%%
\section{\bajoRevision{Fase 3: Análisis de Red y Vulnerabilidad}}
\label{sec:resultados-fase3}

\notaRevision{Los tres indicadores de esta fase —Tiempo de Viaje Promedio,
Centralidad de Intermediación y Vulnerabilidad Estructural— requieren que las
Fases~1 y~2 estén completas. A la fecha de cierre del documento, los algoritmos
de caminos mínimos no han sido ejecutados sobre la versión final ponderada del
grafo. Las fórmulas, interpretaciones esperadas y estructura de resultados
anticipados se presentan como fundamento para la etapa de cálculo.}

La Fase~3 representa la capa de análisis de mayor complejidad computacional del
Modelo VFT. A diferencia de los indicadores de las fases anteriores, que operan
sobre la geometría estática de la red o sobre los pesos individuales de aristas y
nodos, los tres indicadores de esta fase requieren la ejecución de \textbf{algoritmos
de caminos mínimos} sobre la totalidad del grafo ponderado, lo que implica tiempos
de cómputo del orden de horas para un grafo de 11{,}115 nodos y 45{,}679 aristas.

La secuencia de dependencias es estricta: el \textbf{Tiempo de Viaje Promedio}
($T$) exige el grafo con todos sus pesos dinámicos de Fase~2; la
\textbf{Centralidad de Intermediación} ($B$) requiere los caminos mínimos
calculados en $T$; y el análisis de \textbf{Vulnerabilidad Estructural}
($\Delta E$) necesita $B$ para identificar los nodos críticos sobre los cuales
simular fallas.

% ----------------------------------------------------------------
\subsection{\bajoRevision{Tiempo de Viaje Promedio ($T$)}}
\label{subsec:tiempo-viaje-promedio}
% ----------------------------------------------------------------

\notaRevision{Pendiente de cálculo. Requiere Fases 1 y 2 completas más
ejecución del algoritmo \dijkstra{} sobre el grafo final ponderado.}

El \textbf{Tiempo de Viaje Promedio} ($T$) mide la eficiencia global de la red
como sistema: cuánto tarda en promedio un pasajero en desplazarse entre cualquier
par de estaciones alcanzables. Es la métrica de referencia que sintetiza, en un
solo valor, el efecto combinado de la geometría de la red, los tiempos de traslado
con fricción y los costos de transbordo.

\subsubsection{Fórmula}
\label{subsubsec:formula-t}

\begin{equation}
    T = \frac{1}{N(N-1)} \sum_{\substack{i,j \in \mathcal{V} \\ i \neq j}} t(i,j)
    \label{eq:tiempo-viaje-promedio}
\end{equation}

\begin{itemize}
    \item $\mathcal{V}$: conjunto de los $N = 11{,}115$ nodos topológicos del
    grafo.
    \item $t(i,j)$: tiempo de viaje mínimo del nodo $i$ al nodo $j$ en minutos,
    calculado mediante el algoritmo \dijkstra{} sobre el grafo ponderado con los
    atributos de Fase~2.
    \item El denominador $N(N-1)$ normaliza sobre todos los pares ordenados de
    nodos alcanzables. Los pares sin ruta (componentes desconectadas) se excluyen
    del promedio.
\end{itemize}

\subsubsection{Resultado Esperado e Interpretación}
\label{subsubsec:t-resultado-esperado}

El indicador $T$ caracteriza la red en su conjunto como sistema físico: un valor
bajo indica una red eficiente donde cualquier par de estaciones está bien
conectado en el plano temporal; un valor alto sugiere desconexiones estructurales
o cuellos de botella en los transbordos. La comparación de $T$ antes y después de
la remoción simulada de nodos críticos (Sección~\ref{subsec:vulnerabilidad}) es
el mecanismo central del análisis de vulnerabilidad.

Adicionalmente, el modelo contempla una variante ponderada por demanda
$T_{\text{dem}}$, calculable una vez que el \termino{Node Score} de Fase~2 esté
disponible, donde los pares O-D con mayor afluencia en ambos extremos contribuyen
más al promedio. La diferencia $T_{\text{dem}} - T$ cuantificaría el sesgo de la
red: si las rutas más demandadas son sistemáticamente más lentas que el promedio
estructural, ese diferencial tiene implicaciones directas para la política de
mejora de infraestructura.

% ----------------------------------------------------------------
\subsection{\bajoRevision{Centralidad de Intermediación ($B$)}}
\label{subsec:res-centralidad-intermediacion}
% ----------------------------------------------------------------

\notaRevision{Pendiente de cálculo. Requiere la matriz de caminos mínimos
producida por el indicador $T$ (Sección~\ref{subsec:tiempo-viaje-promedio}).}

La \textbf{Centralidad de Intermediación} ($B$) mide cuántas veces cada nodo
aparece en los caminos más cortos entre todos los pares de estaciones de la red.
Un nodo con $B$ alto es un \termino{hub} de paso: su eliminación o degradación
afecta desproporcionadamente al número de rutas disponibles para los pasajeros.
En la hipótesis central de esta investigación, los corredores anulares
—\textit{Periférico}, Circuito Interior— concentran nodos de alta intermediación
que articulan los flujos radiales de la red sin ser ellos mismos estaciones
terminales de alto volumen.

\subsubsection{Fórmula}
\label{subsubsec:formula-b}

\begin{equation}
    B(v) = \frac{1}{(N-1)(N-2)} \sum_{\substack{s,t \in \mathcal{V} \\ s \neq v \neq t}}
    \frac{\sigma_{st}(v)}{\sigma_{st}}
    \label{eq:centralidad-intermediacion}
\end{equation}

\begin{itemize}
    \item $\sigma_{st}$: número total de caminos más cortos (en tiempo) entre
    los nodos $s$ y $t$ en el grafo ponderado.
    \item $\sigma_{st}(v)$: número de esos caminos que pasan por el nodo $v$.
    \item El factor de normalización $1/[(N-1)(N-2)]$ escala el valor al rango
    $[0, 1]$ para un grafo de $N$ nodos.
\end{itemize}

El cálculo se implementa mediante la función \texttt{betweenness\_centrality()}
de \textsc{NetworkX} con el atributo \texttt{weight} como peso de las aristas,
de modo que los caminos mínimos se definen en tiempo, no en número de saltos.

\subsubsection{Resultado Esperado e Interpretación}
\label{subsubsec:b-resultado-esperado}

Se espera que la distribución de $B$ sobre los 11{,}115 nodos sea
\termino{heavy-tailed}: un pequeño número de nodos concentrará valores altos de
intermediación (nodos de paso obligado), mientras que la mayoría de estaciones
tendrá valores cercanos a cero. Este patrón es consistente con la distribución
de red libre de escala (\termino{scale-free}) detectada en la Fuerza Capilar
(Fase~1).

La hipótesis de investigación predice que al menos algunos de los nodos de
mayor $B$ se ubicarán sobre ejes circulares o de transferencia que no son
estaciones de origen-destino frecuente, sino \termino{hubs} de paso que articulan
rutas radiales. Identificar estos nodos es el prerrequisito para el análisis de
vulnerabilidad de la Sección~\ref{subsec:vulnerabilidad}.

% ----------------------------------------------------------------
\subsection{\bajoRevision{Vulnerabilidad Estructural ($\Delta E$)}}
\label{subsec:vulnerabilidad}
% ----------------------------------------------------------------

\notaRevision{Pendiente de cálculo. Requiere la Centralidad de Intermediación
$B$ (Sección~\ref{subsec:res-centralidad-intermediacion}) para seleccionar los nodos
críticos sobre los cuales ejecutar las simulaciones de falla.}

El análisis de \textbf{Vulnerabilidad Estructural} ($\Delta E$) cuantifica en
qué medida la red pierde eficiencia global cuando uno o varios nodos críticos
son removidos. El indicador mide el impacto de fallas selectivas, no aleatorias:
los nodos candidatos a la remoción son aquellos con mayor centralidad de
intermediación $B$, precisamente porque son los que mayor impacto tienen sobre
la conectividad del sistema.

\subsubsection{Fórmula}
\label{subsubsec:formula-deltae}

\begin{equation}
    \Delta E = \frac{E_0 - E_f}{E_0}
    \label{eq:vulnerabilidad}
\end{equation}

\begin{itemize}
    \item $E_0$: eficiencia global de la red original, definida como el inverso
    del tiempo de viaje promedio: $E_0 = 1/T$.
    \item $E_f$: eficiencia global del grafo residual tras la remoción del nodo
    (o conjunto de nodos) crítico.
    \item $\Delta E \in [0, 1]$: un valor cercano a 1 indica que la remoción
    del nodo colapsa la eficiencia de la red; un valor cercano a 0 indica que la
    red es robusta ante esa falla específica.
\end{itemize}

\subsubsection{Protocolo de Simulación}
\label{subsubsec:protocolo-simulacion}

El análisis se ejecutará en tres escalas de falla:

\begin{enumerate}
    \item \textbf{Falla puntual:} remoción de un único nodo (el de mayor $B$)
    y recálculo de $T_f$. Establece el impacto del nodo más crítico individual.
    \item \textbf{Falla por corredor:} remoción simultánea de todos los nodos
    de un eje identificado (e.g., todos los nodos del corredor del Periférico
    Oriente). Mide la vulnerabilidad sistémica del eje completo.
    \item \textbf{Falla en cascada:} remoción secuencial de los $k$ nodos de
    mayor $B$ residual (recalculado tras cada remoción). Simula el escenario de
    falla progresiva en eventos de alta disrupción.
\end{enumerate}

\subsubsection{Resultado Esperado e Interpretación}
\label{subsubsec:deltae-resultado-esperado}

La hipótesis central de esta investigación plantea que la red de transporte de
la \zmvm presenta \textbf{vulnerabilidad asimétrica}: nodos periféricos de alta
intermediación —ubicados en corredores anulares donde convergen rutas radiales—
producen un $\Delta E$ desproporcionadamente mayor al esperado por su posición
geográfica, en comparación con estaciones terminales de alta afluencia como
Pantitlán o Indios Verdes.

Si la simulación confirma esta hipótesis, el análisis proporciona una base
cuantitativa para priorizar inversión en infraestructura de respaldo (rutas
alternativas, capacidad de transbordo) en nodos de alta intermediación pero baja
visibilidad operativa. Si la refuta, el resultado es igualmente valioso: indicaría
que la red posee redundancia suficiente en sus ejes circulares, y que la
vulnerabilidad real se concentra en las terminales radiales de mayor afluencia.

